% !TeX program = pdflatex
\documentclass[11pt]{preprint}
\usepackage[T1]{fontenc}
\usepackage[utf8]{inputenc}
\usepackage{amsmath,amssymb,amsthm}
\usepackage{mathtools}
\usepackage{mathrsfs}
\usepackage[a4paper,margin=28mm]{geometry}
\usepackage{enumitem}
\usepackage{graphicx}
\usepackage{xr-hyper}
\usepackage[hidelinks]{hyperref}
\externaldocument[P-]{WILSON_LOOPS_CONTINUUM_LIMIT}
\newcommand{\dd}{\mathrm{d}}
\newcommand{\tr}{\operatorname{tr}}
\newcommand{\id}{\mathbf 1}
\newcommand{\diam}{\operatorname{diam}}
\newcommand{\dist}{\operatorname{dist}}
\DeclarePairedDelimiter{\abs}{\lvert}{\rvert}
\DeclarePairedDelimiter{\norm}{\lVert}{\rVert}
\DeclarePairedDelimiter{\floor}{\lfloor}{\rfloor}
\DeclarePairedDelimiter{\ceil}{\lceil}{\rceil}
\setlength{\parskip}{0.6em}
\emergencystretch=3em
\title{The Continuum Limit of the Correlation Functions of Wilson Loops in Four-Dimensional Lattice Yang--Mills Theory: Supplement on Ba{\l}aban's Papers}
\author{}
\date{}

\numberwithin{equation}{section}
\newtheorem{proposition}{Proposition}
\renewcommand{\theproposition}{S.\arabic{proposition}}
\newtheorem{lemma}{Lemma}[section]
\theoremstyle{definition}
\newtheorem{definition}[lemma]{Definition}
\theoremstyle{remark}
\newtheorem{remark}[lemma]{Remark}
\theoremstyle{plain}
\renewcommand{\qedsymbol}{\textit{q.e.d.}}
\begin{document}
\maketitle

\begin{abstract}
This supplement contains the statements of Ba{\l}aban's papers on the renormalisation group for four-dimensional lattice gauge theories which the paper uses outside his printed hypotheses, or with a proof which he did not print, each with its proof. A result which Ba{\l}aban proved in his own setting is cited in the paper from his paper and does not appear here.
\end{abstract}

\setcounter{tocdepth}{1}
{\setlength{\parskip}{0pt}\tableofcontents}

\section{Introduction}\label{sec:introduction}

The paper runs Ba{\l}aban's construction with three changes: its parameters are evaluated at numbers fixed by the target instead of the running couplings, its test functions are continuous, and the depth of its $\mathbf R$-operation is chosen by a formula of its own. It inserts a Wilson loop ensemble into the construction, and it compares two runs which differ by one step on the finest lattice. Each of these changes moves the use of some of Ba{\l}aban's results outside the hypotheses under which he proved them, and each of the comparisons needs a result about his objects which he did not state. This supplement collects those results. They are numbered as in the paper, which cites them as ``Supplement, Proposition~S.$n$'', and they are grouped by the part of the paper which uses them.

We use the notation of the paper throughout, and we refer to its statements as ``Paper, Lemma~3.2'' or ``Paper, (3.4)''. Ba{\l}aban's papers are cited by page and equation, as in the paper, with the abbreviations of its bibliography: \cite{BalabanRG1} and \cite{BalabanRG2} for the small-field renormalisation group, \cite{BalabanCRE} for the inductive description of the densities, \cite{BalabanLFR1} and \cite{BalabanLFR2} for the $\mathbf R$-operation and the large-field factors, \cite{BalabanPBG} for the propagators, \cite{BalabanVAR} for the variational problems, \cite{BalabanAVG} for the averaging operations, and \cite{BalabanPRT1,BalabanPRT2} for the random walk expansions. Every constant is one of the paper's, with its place in the order of choices of Paper, Appendix~\ref{P-app:constants}.

\section{The Construction of the Paper}\label{sec:construction}

The construction of Paper, Definition~\ref{P-def:construction}, is Ba{\l}aban's with the parameters of a tube, with soft test functions, and with the depth of Paper, Definition~\ref{P-def:depth}. Paper, Proposition~\ref{P-prop:construction-admissible}, asserts that Ba{\l}aban's description of the densities and his bounds hold for it, and its proof consists of the three propositions of this section, one for each change.

\subsection{The Parameters}\label{sec:parameters}

In \cite{BalabanCRE,BalabanLFR1,BalabanLFR2} every parameter of the scale $j$ is a function of the running coupling $g_j$: the cube size $R_j$, the thresholds $\varepsilon_j = g_jp_0(g_j)$, $\delta_j = (A_1/A_0)\varepsilon_j$, $\delta'_j = g_jA_1p_1(g_j)$, the radii $\alpha_{0,j}$, $\alpha_{1,j}$ of \cite[(2.28), p.~259]{BalabanCRE}, and every threshold, radius, and support of a cutoff function derived from them \cite[(2.4), (2.5), p.~255]{BalabanCRE}, \cite[p.~183]{BalabanLFR1}. In the construction of the paper each of them is the same function evaluated at $\hat g_j$ (Paper, Definitions~\ref{P-def:tube-parameters} and~\ref{P-def:construction}), while the couplings which enter the densities are those of a history of the tube. For the coupling flow of Paper, Section~\ref{P-sec:coupling-flow}, the construction is also run along histories which are not flows, with the factor \eqref{P-eq:history-factor} of the paper at every step, and then the Wilson action of a term of the density of the step $k$ has the coefficient function \eqref{P-eq:history-coupling} of the paper,
\begin{equation}\label{eq:coefficient-function}
    y_k(x) \coloneqq x_k + \sum_{i=1}^{k}\beta_i\bigl(1 - \phi_i(x)\bigr)\;,
\end{equation}
where $\beta_i = \beta_i(x_0, \dots, x_{i-1})$, and $\phi_i \in C_0^\infty(\Lambda_i)$ equals one on $\Lambda_i^{\sim-1}$ \cite[(2.24), p.~259]{BalabanCRE}. Ba{\l}aban does not restrict the functions $\phi_i$ further, and we take them with values in $[0, 1]$. Along a flow, $y_k$ is his coefficient $1/g_k^2(x)$, the solution of the recursion (2.24) which starts from $1/g_0^2(x) = 1/g_0^2$. We put
\begin{equation}\label{eq:action-ratio}
    \bar\rho \coloneqq \frac{15}{8}\;,
\end{equation}
which exceeds the ratio $\rho$ of every tube, by Paper, \eqref{P-eq:tube-ratio-bound}, and we write $C'_W$ for the constant of Paper, Theorem~\ref{P-thm:one-loop-window}. The \emph{tube of a target} is a tube whose constants are those of Paper, Theorem~\ref{P-thm:coupling-flow}: $b = 2c\ln L - \delta'$, $b' = 2c\ln L + \delta'$, and $C_W = C'_W$.

When the Wilson action gives a small factor on a field which has failed a test function, the coefficient of the action enters through a lower bound by the inverse square $g_m^{-2}$ of the running coupling of the scale $m$ of that field. The first step is the corresponding bound for the construction, at every point at which a field of the scale $m$ can be an integration variable.

\begin{lemma}[The Coefficient of the Action]\label{lem:action-coefficient}
Let $\gamma$ satisfy Paper, \eqref{P-eq:gamma-conditions}, the first condition of Paper, \eqref{P-eq:flow-gamma}, and $C'_W\gamma^2 \leq 11/40$. Let a tube be given and let $h = (x_0, \dots, x_k) \in H_k$. Consider a term of the density of the step $k$ along $h$, an integer $m \in \{ 0, \dots, k \}$, and a point $x$ outside $\Lambda_{m+1}$, where $\Lambda_{k+1} \coloneqq \emptyset$.
\begin{enumerate}[label=(\roman*)]
\item If $h$ is a flow, then $y_k(x) \geq \hat g_m^{-2}/\rho$.
\item If $\abs{\beta_i(x_0, \dots, x_{i-1}) - b_i} \leq \delta'$ for $1 \leq i \leq k$, with the number $\delta'$ of Paper, Proposition~\ref{P-prop:coefficients}(i), then $y_k(x) \geq \hat g_m^{-2}/\bar\rho$, and $y_k(x) \geq \hat g_m^{-2}/\rho$ if the tube is the tube of a target.
\end{enumerate}
\end{lemma}

The hypothesis of (ii) is Paper, Proposition~\ref{P-prop:coefficients}(i), at the steps before the step $k$, so that the lemma is available at the step $k$ once the densities up to that step have been constructed. The points outside $\Lambda_{m+1}$ carry the fields of the scale $m$: the block field $V_m$ lives on $\Gamma_m \subset \Omega_{m+1}^c$, and the fluctuation field $A_m$ on $\Omega_{m+1} \setminus \Lambda_{m+1}$ (Paper, \eqref{P-eq:determining-set} and \eqref{P-eq:large-field-operations}). For $m = k$ the lemma holds at every point.

\begin{proof}
We first show that the domains of a term are nested as
\begin{equation}\label{eq:domain-nesting}
    \Lambda_{i+1} \subset \Lambda_i^{\sim-1}\;,
    \qquad
    1 \leq i < k\;,
\end{equation}
or, in terms of the large-field regions, $Z_i^{\sim} \subset Z_{i+1}$. Consider first a term which the step $i + 1$ has produced from a term of the step $i$. Let $Z'_i$ be the union of the cubes of side $LMR_{i+1}$ which meet $Z_i$, the operation $\sim$ being taken for these cubes \cite[p.~264]{BalabanCRE}. By \cite[(3.5), p.~265]{BalabanCRE} the complement of $\Omega_{i+1}$ contains $(B^{i+1}(P^1_{i+1})^c)^{\sim3}$, and $B^{i+1}(P^1_{i+1})^c$ contains $Z_i'^{\sim5}$ by its definition on the same page, so that $\Omega_{i+1} \cap Z_i'^{\sim8} = \emptyset$. The partition into the cubes of side $LMR_{i+1}$ is compatible with the earlier ones \cite[p.~264]{BalabanCRE}, and $LMR_{i+1} \geq MR_i$ by Paper, Lemma~\ref{P-lem:tube-cubes}(ii). Let $x \in \Omega_{i+1}$, let $\square$ be a cube of side $MR_i$ which contains $x$, and let $\square'$ be the cube of side $LMR_{i+1}$ which contains $\square$. Since $\square'$ is not in $Z_i'^{\sim}$, neither $\square'$ nor any cube of side $LMR_{i+1}$ which has a point in common with it meets $Z_i$. Every cube of side $MR_i$ which has a point in common with $\square$ lies in one of these cubes, and therefore in $\Lambda_i$. As $Z_i$ is a union of cubes of side $MR_i$, no point of $Z_i$ lies in $\square$ or in a cube which touches it, that is, $x \in \Lambda_i^{\sim-1}$. Since $\Lambda_{i+1} \subset \Omega_{i+1}$ \cite[(2.1), p.~254]{BalabanCRE}, this proves \eqref{eq:domain-nesting} for the domains which a step produces.

The $\mathbf R$-operation of the step $k$ acts on a union $Z$ of components $X$ of $Z_k$, and the operations on the large-field regions factorise as $\mathbb T_k(Z_k) = \mathbb T_k(Z_k \cap Z^c)\prod_i\mathbb T_k(X_i)$ \cite[(1.1), p.~177]{BalabanLFR1}. A component $X$ which the operation removes is removed together with $\mathbb T_k(X)$, which carries the regions $Z_j \cap X$, $j < k$, by Paper, \eqref{P-eq:large-field-operations}, and the new term is the term $\rho(Z'', V)$ with the remaining large-field regions \cite[(0.3), p.~176]{BalabanLFR1}. This preserves the inclusions $Z_i^{\sim} \subset Z_{i+1}$. Indeed, let $Y$ be a component of $Z_i$ which does not lie in $X$. Then $Y^{\sim}$ is a connected subset of $Z_{i+1} \subset Z_k$ which contains a point outside $X$, so it lies in a component of $Z_k$ other than $X$, and $Y^{\sim} \subset Z_{i+1} \setminus X$.

Since $\phi_{i+1}$ vanishes outside $\Lambda_{i+1} \subset \Lambda_i^{\sim-1}$, while $\phi_i = 1$ on $\Lambda_i^{\sim-1}$ and both take values in $[0, 1]$, we have $\phi_{i+1} \leq \phi_i$ at every point. The weights $w_i \coloneqq 1 - \phi_i(x)$ therefore do not decrease in $i$, and we put $w_0 \coloneqq 0$.

We now sum by parts. For $i > m$ we have $\Lambda_i \subset \Lambda_{m+1}$ by Paper, \eqref{P-eq:domain-sequences}, so that $x \notin \Lambda_i$ and $w_i = 1$. With $S_i \coloneqq \sum_{l=i}^{k}\beta_l$ for $1 \leq i \leq k$ and $S_{k+1} \coloneqq 0$,
\begin{equation}\label{eq:abel-summation}
    y_k(x) = x_k + S_{m+1} + \sum_{i=1}^{m}\beta_iw_i = (1 - w_m)(x_k + S_{m+1}) + \sum_{i=1}^{m}(w_i - w_{i-1})(x_k + S_i)\;.
\end{equation}
The first identity is \eqref{eq:coefficient-function} with $w_i = 1$ for $i > m$. The second follows from the summation by parts $\sum_{i=1}^{m}\beta_iw_i = \sum_{i=1}^{m}(w_i - w_{i-1})(S_i - S_{m+1})$, which holds because $\beta_i = (S_i - S_{m+1}) - (S_{i+1} - S_{m+1})$ and $w_0 = 0$, together with $\sum_{i=1}^{m}(w_i - w_{i-1}) = w_m$. The coefficients $1 - w_m$ and $w_i - w_{i-1}$ are nonnegative and add up to one, so that $y_k(x)$ is a convex combination of the numbers $x_k + S_i$, $1 \leq i \leq m + 1$. Since $u_{i-1} \geq u_m = \hat g_m^{-2}$ for $i - 1 \leq m$, it suffices to show that
\begin{equation}\label{eq:top-down}
    x_k + S_i \geq \frac{u_{i-1}}{\bar\rho}\;,
    \qquad
    1 \leq i \leq k + 1\;,
\end{equation}
under the hypothesis of (ii), and the same with $\rho$ in place of $\bar\rho$ in the other two cases.

(i) If $h$ is a flow, then $x_{l-1} = x_l + \beta_l$ for every $l \leq k$, so that $x_k + S_i = x_{i-1}$, and $x_{i-1} \geq \ell_{i-1} \geq u_{i-1}/\rho$ by Paper, \eqref{P-eq:tube-mediant}.

(ii) For $i = k + 1$ the inequality \eqref{eq:top-down} is $x_k \geq \ell_k \geq u_k/\rho$, by Paper, \eqref{P-eq:tube-mediant}, and $\rho \leq \bar\rho$. For $i \leq k$ we put $n \coloneqq i - 1 < k$. By hypothesis each $\beta_l$ differs from $b_l$ by at most $\delta' \leq \frac25c\ln L$, the first condition of Paper, \eqref{P-eq:flow-gamma}, and $\sum_{l=n+1}^{k}b_l \geq 2c(k - n)\ln L - C'_W$ by Paper, Theorem~\ref{P-thm:one-loop-window}. With $x_k \geq \ell_k = t + b(K - k) - C_W$ this gives
\begin{equation}\label{eq:top-down-lower}
    x_k + S_i \geq (t - C_W - C'_W) + b(K - k) + \tfrac85c\ln L\,(k - n)\;,
\end{equation}
whereas $u_n = (t + C_W) + b'(K - k) + b'(k - n)$. We compare the corresponding terms. Since $C_W \leq \beta_0t/4 < t/8$ by Paper, \eqref{P-eq:tube-class}, and $t \geq \gamma^{-2}$,
\begin{equation}\label{eq:constant-terms}
    \frac{t + C_W}{t - C_W - C'_W} \leq \frac{9/8}{7/8 - C'_W\gamma^2} \leq \frac{15}{8}\;,
\end{equation}
the last step being $C'_W\gamma^2 \leq 11/40$. Moreover $b'/b \leq \frac32$ by Paper, \eqref{P-eq:tube-class}, and $b'/(\frac85c\ln L) \leq \frac{15}{8}$ because $b' \leq b^* = 3c\ln L$ by Paper, \eqref{P-eq:tube-class-constants}. A quotient of two sums of positive numbers is at most the largest quotient of corresponding terms, as in Paper, \eqref{P-eq:tube-mediant}, so that $u_n \leq \bar\rho(x_k + S_i)$, which is \eqref{eq:top-down}.

If the tube is the tube of a target, then $C_W = C'_W$ and $b = 2c\ln L - \delta'$, so that the same two bounds give $S_i \geq b(k - n) - C_W$ and $x_k + S_i \geq (t - 2C_W) + b(K - n)$, while $u_n = (t + C_W) + b'(K - n)$. The quotient of the two is at most $\max((t + C_W)/(t - 2C_W),\, b'/b) = \rho$, by the definition \eqref{P-eq:tube-ratio} of the paper.
\end{proof}

We can now state the proposition. In the displays of \cite[pp.~380--383]{BalabanLFR2} which it concerns, the symbol $p_0(g)$ stands for $(\log g^{-2})^{p_0}$, the factor $A_0$ of \cite[(2.4), p.~255]{BalabanCRE} being written out, as in $p_1(g) = (\log g^{-2})^{p_1}$ \cite[p.~183]{BalabanLFR1}. Only in this reading does the printed bound of $\exp(-\frac14g_m^{-2}B_3^{-2}\frac12\varepsilon_m^2)$ by $\exp(-\frac18B_3^{-2}(1 + \beta_0)^{-2}A_0^2p_0^2(g_j))$, for $j - N \leq m \leq j$ \cite[p.~381]{BalabanLFR2}, follow from $g_m^{-2}\varepsilon_m^2 = A_0^2(\log g_m^{-2})^{2p_0}$ and (2.7), and we read all of them so (Paper, Section~\ref{P-sec:notes-on-papers}).

\begin{proposition}[The Parameters of a Tube]\label{prop:parameters}
Let $\gamma$ satisfy the hypotheses of Lemma~\ref{lem:action-coefficient}, let $\beta_0 \leq \sqrt2 - 1$, and let $A_1$ and $A_0$ satisfy
\begin{equation}\label{eq:threshold-constants}
    A_1 \geq 48\bar\rho^{1/2}\;,
    \qquad
    A_0^2 \geq \bar\rho A_1^2\max\bigl(8(1 + \beta_0)^2B_3^2,\ 4O(1)B_3M^2\bigr)\;,
\end{equation}
with the constant $O(1)$ of \cite[p.~381]{BalabanLFR2}. Let a tube be given, and let $h \in H_k$ be a history which is a flow or whose coefficients satisfy the hypothesis of Lemma~\ref{lem:action-coefficient}(ii). Every statement of \cite{BalabanCRE,BalabanLFR1,BalabanLFR2}, and of the parts of the papers which they import, in which a parameter of a scale enters holds for the construction of the paper along $h$, with the following changes:
\begin{enumerate}[label=(\alph*)]
\item a relation among the parameters alone holds as printed;
\item a lower bound drawn from the action, in which the Wilson action gives a factor $\exp(-cF(\log g_j^{-2}))$ on a field above a threshold, with $F$ nondecreasing, holds with $c$ replaced by $c/\bar\rho$, and by $c/\rho$ along a flow and on the tube of a target;
\item an upper bound relative to the action holds with the printed constant, or with a constant which differs from it by a factor at most $2$;
\item a conversion between a scaled and an unscaled field holds with the coupling of the history as the factor and the parameter of the tube as the threshold;
\item the $\mathbf R$-terms produced by the step $k$ satisfy \cite[(2.31), p.~260]{BalabanCRE} with the coupling $g_k$ of the history.
\end{enumerate}
In particular, given Propositions~\ref{prop:soft-tests} and~\ref{prop:depth} for the test functions and the depth, the conclusions of Paper, Theorem~\ref{P-thm:balaban-inductive}, and the bounds \eqref{P-eq:balaban-theorem-two} and \eqref{P-eq:ultraviolet-stability} of the paper hold for the densities of the construction along $h$, with constants which do not depend on the tube, on $h$, on $k$ and $K$, or on the torus.
\end{proposition}

The conditions \eqref{eq:threshold-constants} are of the kind which Ba{\l}aban imposes himself, a lower bound on $A_1$ and then on $A_0$ in terms of earlier constants \cite[p.~246, p.~265]{BalabanCRE}, and they belong to the stage $6$ of Paper, Appendix~\ref{P-app:constants}, as does $C'_W\gamma^2 \leq 11/40$.

\begin{proof}
Theorem~1 of \cite[p.~262]{BalabanCRE} assumes the inequalities (0.31) of \cite[p.~259]{BalabanRG1}, which it cites as (I.0.33) (Paper, Section~\ref{P-sec:notes-on-papers}), and its proof uses them through the inequalities (2.6)--(2.9) of \cite[pp.~255--256]{BalabanCRE}, which Ba{\l}aban derives from the renormalisation group equations and uses in their place \cite[p.~255]{BalabanCRE}. A history of the tube satisfies (2.6) and (2.7) by Paper, Lemma~\ref{P-lem:tube-couplings}(ii), with real exponents, and so does the sequence $(\hat g_j)$, by part (iii) of that lemma; the parameters of the tube satisfy (2.8) by Paper, Lemma~\ref{P-lem:tube-thresholds}(ii), and the cube sizes satisfy (2.5) and (2.9) by Paper, Lemma~\ref{P-lem:tube-cubes}. It remains to go through the places at which a parameter enters, and we sort them into the kinds (a)--(e) and two further kinds.

\emph{Relations among the parameters.} These are the compatibility of the partitions and the choice of the cube sizes \cite[pp.~245, 264]{BalabanCRE}, \cite[p.~177]{BalabanLFR1}; the definitions of $k_1$ \cite[p.~270]{BalabanCRE} and of $N_0$ \cite[p.~179]{BalabanLFR1}; the inequalities (2.6)--(2.9) and their consequences, such as $L^{N_0-1} = R_{k-N_0+1} \leq (L + 1)(N_0 + 1)^{\beta_0}R_k$ and $\delta'_j \leq (1 + \beta_0)^2(k - j)^{\beta_0}\delta'_k$ \cite[p.~192]{BalabanLFR1}; and the margins between powers of logarithms, such as that of $2(1 + \beta_0)^{-1}p_0(g_{j+1})$ against $O(1)3(100M)^dR_{j+1}^{d+2}$ \cite[(1.88), p.~387]{BalabanLFR2}. In the construction all of them are evaluated on the sequence $(\hat g_j)$, which lies in $(0, \gamma]$ and satisfies Ba{\l}aban's hypotheses, as we have just seen. So they hold as printed.

\emph{Lower bounds from the action.} Here the Wilson action gives a small factor on a field of the scale $m$ which exceeds a threshold. Ba{\l}aban bounds the coefficient of the action below by $g_m^{-2}$, the threshold being $\tau(g_n)$ for a scale $n$ and a function $\tau(g) = cg(\log g^{-2})^p$, and then bounds $g_m^{-2}\tau(g_n)^2$ below by a nondecreasing function $F$ of $\log g_j^{-2}$, by means of (2.6)--(2.8). In the construction the coefficient at the points which carry the field is at least $\hat g_m^{-2}/\bar\rho$, by Lemma~\ref{lem:action-coefficient}, and the threshold is $\tau(\hat g_n)$. The same derivation, applied to the sequence $(\hat g_j)$, which satisfies (2.6)--(2.8), therefore gives
\begin{equation}\label{eq:action-comparison}
    y_k\tau(\hat g_n)^2 \geq \frac{1}{\bar\rho}\,\hat g_m^{-2}\tau(\hat g_n)^2 \geq \frac{1}{\bar\rho}F(\log\hat g_j^{-2}) \geq \frac{1}{\bar\rho}F(\log g_j^{-2})\;,
\end{equation}
the last step because $\hat g_j \leq g_j$ and $F$ does not decrease. So the factor is Ba{\l}aban's with the constant in the exponent divided by $\bar\rho$, and by $\rho$ along a flow and on the tube of a target. Examples are the factors $\exp(-p_0(g_0))$ and $\exp(-\frac12p_0^2(g_0))$ of a failed plaquette test function \cite[p.~175]{BalabanLFR1}, \cite[p.~244]{BalabanCRE}, the large-field factors of \cite[pp.~380--385]{BalabanLFR2}, and the Gaussian tails of the fluctuation integrals. A change of the constant in the exponent is harmless wherever the bound has a margin of a positive power of $\log g_j^{-2}$, and it remains to check the places where two such factors are compared at the same order. They are the following, and each of them holds with the factor $\bar\rho^{-1}$ in the exponent on the side of the action.
\begin{enumerate}[label=(\roman*)]
\item The creation factor: \cite[p.~381]{BalabanLFR2} bounds $\exp(-\gamma_0\min\{\frac12B_3^{-2}A_0^2, 2A_1^2, A_1^2\}p_0^2(g_j)(d'_j + 1))$ by $\exp(-\frac12\gamma_0A_1^2p_0^2(g_j)(d'_j + 1) - 2p_0(g_j))$, where $d'_j \geq 0$ is the linear size of a component in terms of the cubes of side $MR_j$ \cite[p.~380]{BalabanLFR2}. By \eqref{eq:threshold-constants} the minimum is $A_1^2$. With $\gamma_0/\bar\rho = \frac12\gamma_0 + \frac{1}{30}\gamma_0$ in place of $\gamma_0$ the left side is at most $\exp(-\frac12\gamma_0A_1^2p_0^2(g_j)(d'_j + 1) - \frac{1}{30}\gamma_0A_1^2p_0^2(g_j))$, because $d'_j + 1 \geq 1$, and $\frac{1}{30}\gamma_0A_1^2p_0^2(g_j) \geq 2p_0(g_j)$ for $g_j \leq \gamma$ small.
\item \cite[p.~381]{BalabanLFR2} bounds $\exp(-\frac18B_3^{-2}(1 + \beta_0)^{-2}A_0^2p_0^2(g_j))$ by $\exp(-A_1^2p_0^2(g_j))$. With $\bar\rho^{-1}$ in the first exponent this is the first part of the second condition of \eqref{eq:threshold-constants}.
\item \cite[p.~381]{BalabanLFR2} bounds $\exp(-\gamma_0g_m^{-2}(2\delta'_m)^2)$ by $\exp(-\gamma_04(1 + \beta_0)^{-2}A_1^2p_1^2(g_j))$ and then by $\exp(-\gamma_0A_1^2p_1^2(g_j))$. With $\bar\rho^{-1}$ the second step needs $4(1 + \beta_0)^{-2} \geq \bar\rho$, which holds because $\beta_0 \leq \sqrt2 - 1$ gives $4(1 + \beta_0)^{-2} \geq 2$.
\item \cite[p.~381]{BalabanLFR2} bounds $\exp(-\frac14g_j^{-2}(O(1)B_3M^2)^{-1}\varepsilon_j^2)$ by $\exp(-A_1^2p_0^2(g_j))$. Since $g_j^{-2}\varepsilon_j^2 = A_0^2p_0^2(g_j)$ in the reading above, with $\bar\rho^{-1}$ this is the second part of the second condition of \eqref{eq:threshold-constants}.
\item \cite[p.~382]{BalabanLFR2} obtains the factor $\exp(-\frac14g_j^{-2}(24R_j^4)^{-2}\delta_j'^2) = \exp(-A_1^248^{-2}R_j^{-8}p_1^2(g_j))$ and bounds it by $\exp(-R_j^{-8}p_1^2(g_j))$. With $\bar\rho^{-1}$ the last step is the first condition of \eqref{eq:threshold-constants}.
\item The bound $\frac14\gamma_0(14)^{-d}A_1^2p_0^2(g_1) \geq O(1)2(64)^dM^dL^{d+1}R_1^{d+2}$ \cite[p.~385]{BalabanLFR2} concerns the exponent obtained from (i), which holds as printed, and it has the margin of the factor $(\log g_1^{-2})^{2p_0}$ against $R_1^{d+2} < (2L)^{d+2}(\log g_1^{-2})^{(d+2)r}$, by Paper, Lemma~\ref{P-lem:tube-cubes}(i), with $2p_0 \geq 10r > (d + 2)r$ at $d = 4$.
\end{enumerate}

\emph{Upper bounds relative to the action.} Here a field below a threshold enters an upper bound, often multiplied by the coefficient of the action, as in the bound $O(1)\abs{\mathbf B_0}\delta_k'^3 \leq g_k^3O(1)M^dR_k^{d+2}A_1^3p_1^3(g_k)$ \cite[p.~361]{BalabanLFR2} and in the bound $O(1)g_k^{-2}(\alpha_{0,k} + \alpha_{1,k})\varepsilon_kL^{-N}N^{1+\beta_0}$ \cite[p.~368]{BalabanLFR2}. The thresholds of a tube are the smaller ones, $\tau_j \leq \tau(g_j)$ by Paper, \eqref{P-eq:tube-threshold-ratio}. The coefficient of the action exceeds $x_k = g_k^{-2}$ by at most the sum of the numbers $\abs{\beta_i}$ over the scales $i$ involved, and each of them is at most $2c\ln L + C'_W + \delta'$, by Paper, Theorem~\ref{P-thm:one-loop-window} with $n = m + 1$ and Proposition~\ref{P-prop:coefficients}(i); this is the form $1/g_k^2 + O(k - j)$ in which the coefficient enters \cite[p.~368]{BalabanLFR2}. So a bound of this kind holds as printed when it uses only upper bounds on the thresholds. Where it uses a quotient of two parameters of the same scale, the quotient at $\hat g_j$ differs from the quotient at $g_j$ by a factor $(\log\hat g_j^{-2}/\log g_j^{-2})^{\pm p}$ with $p \leq p^*$, which lies between $1/2$ and $2$ by Paper, Lemma~\ref{P-lem:tube-couplings}(i), and the third condition of Paper, \eqref{P-eq:gamma-conditions}.

\emph{Conversions between scaled and unscaled fields.} These are the conditions $\abs{A(b)} < g_0^{-1}\delta_0$ \cite[p.~250]{BalabanCRE}, $\abs{A_j(b)} < g_j^{-1}\delta_j$ \cite[(1.9), p.~178]{BalabanLFR1}, the condition on $B_j$ of \cite[(1.55), p.~188]{BalabanLFR1}, and $\abs{B(b)} \geq g_j^{-1}\delta'_j$ \cite[p.~383]{BalabanLFR2}. In the construction the factor $g_j$ is the coupling of the history and the threshold is the parameter of the tube, by Paper, Definition~\ref{P-def:construction}(ii), so that each of them is a condition on the unscaled field with the threshold of the tube. The domains of analyticity, such as the domain (3.43) of \cite[p.~276]{BalabanCRE}, in which $\abs{A} < C_1p_1(g_k)$, are domains on which bounds are asserted and not thresholds under integrals, and they are evaluated at $\hat g_j$ like every other parameter.

\emph{The $\mathbf R$-terms.} The terms which the $\mathbf R$-operation produces at the step $k$ satisfy $\abs{\mathbf R'^{(k)}(X, (\mathbf U, \mathbf J))} \leq \exp(-p_0(\hat g_k))\exp(-\kappa d_k(X))$ \cite[p.~390]{BalabanLFR2}, with the parameter of the tube. Since $p_0(g) = A_0(\log g^{-2})^{p_0} \geq \frac{\kappa_0}{2}\log g^{-2}$ for $g \leq \gamma$ small, $\exp(-p_0(\hat g_k)) \leq \hat g_k^{\kappa_0} \leq g_k^{\kappa_0}$, which is \cite[(2.31), p.~260]{BalabanCRE} with the coupling of the history.

\emph{Quantities of the history and the geometry.} Two further uses of the parameters are not comparisons with the action. Quantities of the history which are bounded by a parameter, such as the bound $O(N) \leq O(R_k)$ on $1/(g''_k(\cdot))^2 - 1/g_k^2$ \cite[p.~366]{BalabanLFR2} and the terms ``bounded by any positive power of $g$'' \cite[pp.~252, 269]{BalabanCRE}, \cite[p.~188]{BalabanLFR1}, have the history on one side, which is bounded through the coefficients $\beta_i$ and the couplings $g_j \geq \hat g_j$, and a parameter of the tube on the other, which differs from Ba{\l}aban's by the factors of Paper, \eqref{P-eq:tube-threshold-ratio}. The domains, the partitions, and the cube sizes are fixed by the target and the configuration, and the combinatorial bounds on the numbers of polymers, cubes, and domains hold as printed with the cube sizes of the tube. The lower bounds on $R_j$, such as the smallness of $\exp(-R_k)$ \cite[pp.~266, 268, 269]{BalabanCRE} and $\log g_j^{-2} \leq R_j$ \cite[p.~380]{BalabanLFR2}, hold by Paper, Lemma~\ref{P-lem:tube-cubes}(i).

\emph{The imported papers.} The papers \cite{BalabanAVG,BalabanVAR,BalabanREG} are statements about lattice gauge fields under a smallness condition on the field, and they contain no coupling; they apply to the thresholds of the tube, which are smaller than Ba{\l}aban's. The propagator and random walk estimates \cite{BalabanPBG,BalabanPRT1,BalabanPRT2} contain no coupling either. The small-field steps of \cite{BalabanRG1,BalabanRG2}, which \cite{BalabanCRE} and \cite{BalabanLFR2} run on their own domains, use the parameters as smallness conditions on thresholds, as the threshold of the fluctuation field relative to the radii $\alpha_{0,j}$, $\alpha_{1,j}$, or in relations among the parameters, so that they fall under the kinds above.

Every place at which a parameter enters the three papers and the parts of the imported papers which they use is of one of the kinds (a)--(e) or of the two further kinds, and in each case the statement holds for the construction with constants which depend on the constants of the construction and not on the tube, on $h$, on $k$ and $K$, or on the torus. Ba{\l}aban's proof therefore reproduces the description of the densities under these changes, given the corresponding statements for the test functions and the depth, which are Propositions~\ref{prop:soft-tests} and~\ref{prop:depth}.
\end{proof}

\subsection{Soft Test Functions}\label{sec:soft-tests}

A soft test function differs from a characteristic function in two ways, and each of them has to be followed through Ba{\l}aban's papers. Its support is larger, so that a configuration on which a passed test function does not vanish obeys the bound $f_\square < (1 + \kappa_s)t$ instead of $f_\square < t$; and it is not idempotent, so that an argument which multiplies a test function by itself, or removes it because it equals one, has to be checked. The first point is a question of margins: wherever Ba{\l}aban draws a bound from a passed test function and then uses the bound to show that another test function equals one, or that a field lies in a domain, his inequality must have room for the factor $1 + \kappa_s$. The second point concerns the cascade of test functions of the $\mathbf R$-operation \cite[(1.22)--(1.24), p.~181]{BalabanLFR1}, in which each test function is removed on the support of the next one, and there the widths have to be nested.

We first say which functions are soft. The cascade of the $\mathbf R$-operation of the step $k$ consists of the test functions $\chi_k^{(n)}$, $n = 0, \dots, N$, where $N = k - h$ is the number of integrations which the operation performs and $h$ the scale at which it starts \cite[pp.~181--182]{BalabanLFR1}. The test function $\chi_k^{(n)}$ bounds the plaquettes of the configuration $U^{(n)}_{k,Z}$ by the thresholds $T_n(p)$ of \cite[(1.24), p.~182]{BalabanLFR1}, which depend on the layer of the plaquette $p$.

\begin{definition}[The Soft Functions of the Construction]\label{def:soft-functions}
In the construction of the paper the following functions are soft, with the width $\kappa_s = \beta/30$ of Paper, \eqref{P-eq:soft-width}, unless stated otherwise.
\begin{enumerate}[label=(\roman*)]
\item The test functions of the plaquettes and of the reference configurations, \cite[(1.1), (1.4), (1.8), (2.17), (3.2), (3.3)]{BalabanCRE}, \cite[(1.3), (1.4), (1.7), (1.8), (1.27), (1.75), (1.88)]{BalabanLFR1}, and their restatements in \cite[(1.70), p.~380]{BalabanLFR2}: each characteristic function $\chi(\{ f_\square < t \})$ is replaced by $\bar\chi(f_\square/t)$, and its complement by $1 - \bar\chi(f_\square/t)$, with $f_\square \coloneqq \infty$ where the configuration which defines it does not exist.
\item The cascade: for a cube $\square$ let $f^{(n)}_\square$ be the supremum of $\abs{U^{(n)}_{k,Z}(\partial p) - 1}/T_n(p)$ over the plaquettes $p$ of $\chi_k^{(n)}$ assigned to $\square$, each plaquette being assigned to one cube of side $LM_2R_i$ which contains it, $i$ the scale of $T_n(p)$. Then $\chi_k^{(n)}$ is replaced by $\prod_\square\bar\chi_c(f^{(n)}_\square/(1 + \kappa_c)^{N-n})$, where $\bar\chi_c$ is the function \eqref{P-eq:soft-function} of the paper with the width
\begin{equation}\label{eq:cascade-width}
    \kappa_c \coloneqq \frac{\kappa_s}{N + 1}
\end{equation}
in place of $\kappa_s$, and its complement accordingly.
\item The cutoff $\chi(\Lambda_i)$ of \cite[(1.101), p.~201]{BalabanLFR1} is replaced by $\prod_b\bar\chi(\abs{(1/i)\log V'(b)}/(M_0\varepsilon_k))$, the product over the bonds $b$ of $\Lambda_i$, with $M_0 \coloneqq 2(1 + \kappa_s)O(1)B_5M^6$, where $O(1)B_5M^6$ is the bound of \cite[p.~378]{BalabanLFR2}.
\end{enumerate}
The characteristic functions of the fluctuation fields and of the axial gauge fixing, \cite[(1.5), (1.18), (3.16)]{BalabanCRE} and \cite[(1.5), (1.9), (1.55), (1.82)]{BalabanLFR1}, are sharp.
\end{definition}

The number $M_0$ is left free by Ba{\l}aban, who sets it equal to the bound $O(1)B_5M^6$ on $\abs{V'_k - 1}$ which the passed test functions give \cite[p.~378]{BalabanLFR2}; the factor $1 + \kappa_s$ accounts for the soft test functions and the factor $2$ for the logarithm, since $\abs{(1/i)\log v} \leq \frac{\pi}{2}\abs{v - 1}$ for a unitary $v$. The nesting of the cascade is the content of the next lemma.

\begin{lemma}[The Soft Cascade]\label{lem:soft-cascade}
The soft test function $\chi_k^{(n)}$ equals one where $f^{(n)}_\square \leq (1 + \kappa_c)^{N-n}$ for every $\square$, and it vanishes unless $f^{(n)}_\square < (1 + \kappa_c)^{N-n+1}$ for every $\square$. Every threshold of a soft cascade test function lies between $T_n(p)$ and $(1 + \kappa_c)^{N+1}T_n(p) \leq (1 + 2\kappa_s)T_n(p)$, and $x \mapsto \bar\chi_c(x/(1 + \kappa_c)^{N-n})$ is Lipschitz with the constant $(N + 1)/\kappa_s$.
\end{lemma}

\begin{proof}
The first two statements are those of $\bar\chi_c$, which equals one on $[0, 1]$ and vanishes on $[1 + \kappa_c, \infty)$, at the argument $x/(1 + \kappa_c)^{N-n}$. For the third, $(1 + \kappa_c)^{N+1} = (1 + \kappa_s/(N + 1))^{N+1} \leq e^{\kappa_s} \leq 1 + 2\kappa_s$, the last step because $e^x \leq 1 + (e - 1)x$ for $0 \leq x \leq 1$, by convexity. The Lipschitz constant of $\bar\chi_c$ is $1/\kappa_c = (N + 1)/\kappa_s$, and the division of the argument by $(1 + \kappa_c)^{N-n} \geq 1$ does not increase it.
\end{proof}

\begin{proposition}[Soft Test Functions]\label{prop:soft-tests}
Let $\beta = 1/8$, $\kappa_s = \beta/30$, $\beta_0 \leq 1/7$, $L \geq 3$, and let $\gamma$ be small. For every tube and every history as in Proposition~\ref{prop:parameters}, the statements of \cite{BalabanCRE,BalabanLFR1,BalabanLFR2} hold for the construction with the soft functions of Definition~\ref{def:soft-functions}, with the following changes.
\begin{enumerate}[label=(\alph*)]
\item The decompositions of unity are exact, and every removal of a test function which Ba{\l}aban proves on the support of the other test functions holds on the plateau of the removed one, given the others at the thresholds of their supports.
\item Every bound drawn from a passed test function holds with the threshold $t$ replaced by $(1 + \kappa_s)t$, and by $(1 + 2\kappa_s)t$ for a cascade test function; every bound drawn from a failed test function holds at the threshold $t$.
\item The changes of variables \cite[(1.22), p.~251]{BalabanCRE}, its repetition at \cite[p.~269]{BalabanCRE}, and \cite[(1.59), p.~188]{BalabanLFR1} carry a soft test function into the soft test function of the new variables.
\item The normalisation \cite[(1.102), p.~201]{BalabanLFR1} holds, and the denominator of \cite[(1.100)]{BalabanLFR1} with the soft cutoff lies within the factors $e^{\pm1/2}$ of the denominator with the sharp cutoff at the threshold $M_0\varepsilon_k$.
\item In the term of the density of the step $k$ which has no large-field region, every test function other than $\chi_k$ of Paper, \eqref{P-eq:test-function}, equals one on the support of $\chi_k$.
\end{enumerate}
\end{proposition}

\begin{proof}
(a) The identity \eqref{P-eq:soft-decomposition} of the paper holds for any functions with values in $[0, 1]$, by Paper, Lemma~\ref{P-lem:soft-decomposition}(i), and so the lattice integrals are unchanged whichever decomposition expands them. Where test functions on overlapping families of cubes are combined, as $\chi_k(X_i^{\sim-6})$ with $\chi_k(\Omega_k^{\sim4})$ \cite[p.~378]{BalabanLFR2}, the soft product keeps both factors, with values in $[0, 1]$ and with the support of each, and no identity $\bar\chi^2 = \bar\chi$ is used. A removal is a statement of the form: on the support of the remaining test functions the quantity $f$ of the removed one satisfies $f < t$. With soft test functions the supports are those of Paper, Lemma~\ref{P-lem:soft-decomposition}(ii), and the conclusion needed is $f \leq t$, the plateau, so that it is a bound drawn from passed test functions and falls under (b), with one exception, the cascade, which uses no margin at all. Indeed, Ba{\l}aban removes $\chi_k^{(n)}$ on the support of $\chi_k^{(n+1)}$ by inequalities of the form $f \leq af' + y$ at each plaquette, with $f$ the quantity of $\chi_k^{(n)}$, $f'$ that of $\chi_k^{(n+1)}$, and numbers $a, y \geq 0$ which do not depend on the threshold of $\chi_k^{(n+1)}$: this is \cite[(1.48), p.~186]{BalabanLFR1}, with $a = 1 + \alpha\beta2^{-(j-i)}$ and $y = \alpha\beta2^{-(j-i)}\varepsilon_i(L^{k-i}\eta)^2$, used in \cite[(1.50)--(1.52), p.~187]{BalabanLFR1}, where he shows $aT_{n+1}(p) + y \leq T_n(p)$. On the support of the soft $\chi_k^{(n+1)}$ we have $f' < (1 + \kappa_c)^{N-n}T_{n+1}(p)$ by Lemma~\ref{lem:soft-cascade}, and therefore
\begin{equation}\label{eq:cascade-removal}
    f \leq (1 + \kappa_c)^{N-n}\bigl(aT_{n+1}(p) + y\bigr) \leq (1 + \kappa_c)^{N-n}T_n(p)\;,
\end{equation}
because $y \leq (1 + \kappa_c)^{N-n}y$. This is the plateau of the soft $\chi_k^{(n)}$, so the removal holds exactly. The coefficients $a$ and $y$ come from numerical factors which Ba{\l}aban makes at most $\alpha\beta$ by taking $\gamma$ or $A_1/A_0$ small \cite[p.~186]{BalabanLFR1}, and the factor $1 + \kappa_s$ of the soft test function \cite[(1.27)]{BalabanLFR1}, which enters through the field $(1/i)\log[V_j(V^{(j)}_Z)^{-1}]$, is absorbed in the same way.

(b) A failed test function enters as a lower bound $f > t$ on its quantity, from which the Wilson action gives a small factor. Since $1 - \bar\chi(f/t)$ vanishes for $f \leq t$, such a bound holds at the threshold $t$, as printed. This applies to the gains of \cite[pp.~380--381]{BalabanLFR2}, which are monotone in the threshold, and to the failed test functions of \cite{BalabanCRE,BalabanLFR1}.

A passed test function enters as an upper bound $f < (1 + \kappa_s)t$, or $f < (1 + 2\kappa_s)t$ for the cascade by Lemma~\ref{lem:soft-cascade}, and its consequences are of two kinds. Most of them carry a constant, a factor $\exp(-\delta MR_j)$ or $\exp(-R_j)$, a factor $L^{-N}$, or a negative power of $\log\hat g_j^{-2}$ such as $A_1/A_0$ or $\delta'_j/\varepsilon_j$, and are then used for $\gamma$ small, against such a factor, or as a regularity condition with an unspecified constant, and the factor $1 + 2\kappa_s$ is absorbed. Examples are \cite[(3.7), p.~266]{BalabanCRE}, where the boundary-layer term $30d^2L^2B_3^2(1 + \beta_0)\exp(-R_k)\varepsilon_k$ and the term $4B_3(A_1/A_0)\varepsilon_k$ scale by the factor; \cite[(3.9), p.~266]{BalabanCRE}, $O(1)(A_1/A_0)\varepsilon_k < \varepsilon_k$; \cite[(3.13), p.~267]{BalabanCRE}, $(4dL)^2(1 + \beta_0)B_3\varepsilon_k$; the regularity $\abs{\partial U_{k+1} - 1} < 2B_3\varepsilon_{k+1}(L^{-1}\eta)^2$ and the bound $B_3\exp(-\delta LM_2R_{k+1})44d^2B_3\varepsilon_{k+1}$ of \cite[p.~268]{BalabanCRE}; $O(1)\delta'_j < \varepsilon_j$ \cite[p.~184]{BalabanLFR1}; $O(1)B_3\exp(-\delta LM_2R_{j+1})22d^2(1 + \beta_0)(A_0/A_1)\delta_j \leq \frac12\delta_j$ \cite[p.~185]{BalabanLFR1}; the regularity $\abs{\partial V_k - 1} < 2L_0^{2N_0}\varepsilon_k$, absorbed in $O(1)M^2L^2R_k\varepsilon_k$ \cite[p.~193]{BalabanLFR1}; and in \cite[pp.~356--377]{BalabanLFR2} the bounds linear or quadratic in $\varepsilon_k$ with a constant, such as $44d^2\varepsilon_k$ \cite[p.~365]{BalabanLFR2}, used for $\gamma$ small or against $L^{-N}$ \cite[p.~363]{BalabanLFR2}.

The others compare a quantity with a threshold at the same order, and there the margin has to be checked. They are the following.
\begin{enumerate}[label=(\roman*)]
\item \cite[(1.31), p.~184]{BalabanLFR1}, the smallest margin. The plaquettes of $U_{j,\square}$ are bounded by $(1 - \beta/2 + (2\beta/10)L^{-2} + 4\beta/10)\varepsilon_j\xi^2 = (1 - \beta/10 + \beta/(5L^2))\varepsilon_j\xi^2$, from the last inequality of (1.24) for the cascade and the bound $B_3\exp(-2\delta M_2R_j)22d^2\varepsilon_j < (\beta/10)\varepsilon_j$ of the boundary term. With the cascade threshold enlarged by $1 + 2\kappa_s = 1 + \beta/15$, and the boundary term by the same factor, which is absorbed for $\gamma$ small, the bound becomes
\begin{equation}\label{eq:smallest-margin}
    \Bigl(1 + \frac{\beta}{15}\Bigr)\Bigl(1 - \frac{\beta}{10} + \frac{\beta}{5L^2}\Bigr) \leq \Bigl(1 + \frac{\beta}{15}\Bigr)\Bigl(1 - \frac{7\beta}{90}\Bigr) = 1 - \frac{\beta}{90} - \frac{7\beta^2}{1350} < 1
\end{equation}
times $\varepsilon_j\xi^2$, where the first step is $L \geq 3$. So the test functions (1.3), and (1.4) for $j + 1 < k$, are on their plateaus, as on \cite[p.~184]{BalabanLFR1}.
\item \cite[(3.8), p.~266]{BalabanCRE}: the bound $(2(1 + \beta_0)L^{-2} + 4/10)\varepsilon_k\eta^2$ becomes $(1 + \kappa_s)(2(1 + \beta_0)L^{-2} + 4/10)\varepsilon_k\eta^2 \leq 0.66\,\varepsilon_k\eta^2$ for $L \geq 3$ and $\beta_0 \leq 1/7$, which is below the plateau $\varepsilon_k\eta^2$.
\item \cite[(1.96), p.~199]{BalabanLFR1}: with $\alpha = 1/12$ \cite[p.~199]{BalabanLFR1}, the bound of the plaquettes of $U''_{k,Z}$ drawn from the passed $\chi_{h,1/2}$ at the threshold $\frac12(1 + \kappa_s)\varepsilon_h$ is $(\frac12(1 + \kappa_s) + 2\alpha + O(\varepsilon_h))\varepsilon_h(L^{k-h}\eta)^2 = (\frac23 + \frac{\kappa_s}{2} + O(\varepsilon_h))\varepsilon_h(L^{k-h}\eta)^2$, which is below Ba{\l}aban's $\frac34\varepsilon_h(L^{k-h}\eta)^2$ for $\gamma$ small, since $\kappa_s < 1/6$.
\item \cite[p.~199]{BalabanLFR1}: the bound $22d^2\varepsilon_h + 2\delta'_k < 23d^2\varepsilon_h$ becomes $22d^2(1 + 2\kappa_s)\varepsilon_h + 2\delta'_k < 23d^2\varepsilon_h$, which holds for $\gamma$ small because $44d^2\kappa_s < d^2$.
\item \cite[(1.98), p.~200]{BalabanLFR1}: at $j = k$ the leading term $2\varepsilon_k\eta^2$, drawn from the passed $\chi_{k,\Lambda}$, becomes $2(1 + \kappa_s)\varepsilon_k\eta^2$, and the bound $(2(1 + \kappa_s) + 5\alpha + O(\varepsilon_k))\varepsilon_k\eta^2 \leq 2.43\,\varepsilon_k\eta^2$ lies below the threshold $3(1 - \beta/2)\varepsilon_k\eta^2 = 2.8125\,\varepsilon_k\eta^2$ of (1.24) at $c = 3$; at $j < k$ the bound has the room of the factor $L^{-(k-j)}$.
\end{enumerate}

(c) The change of variables \cite[(1.22), p.~251]{BalabanCRE} contains a smooth function $g$ with $g = 1$ on a neighbourhood of the set where the tested quantity of the reference test function reaches the threshold $2\delta_0$ \cite[p.~252]{BalabanCRE}. A soft test function and its complement differ from $1$ and from $0$ only where the tested quantity is at least the threshold, and there $g = 1$, so that the two arguments coincide; where $g < 1$ the tested quantity is below the threshold, and both the soft and the sharp test function equal one. The change of variables therefore carries the soft test function into the soft test function of the new variables. The same holds at \cite[p.~269]{BalabanCRE} and for \cite[(1.59), p.~188]{BalabanLFR1}, which are changes of variables of the same form.

(d) The cutoff of Definition~\ref{def:soft-functions}(iii) is the same function in the numerator of \cite[(1.100)]{BalabanLFR1}, at $V' = V'_k$, and in its denominator, an integral over $V'$. At fixed $V_k$ outside $\Lambda$ the function $V_\Lambda$ does not depend on $V_k\lceil_\Lambda$, so that $V_k\lceil_\Lambda \mapsto V'_k = V_k(V_\Lambda)^{-1}$ is a translation of the Haar measure, and the integral of the quotient over $V_k\lceil_\Lambda$ is one, which gives (1.102). Since $\bar\chi(y) = \kappa_s^{-1}\int_1^{1+\kappa_s}\mathbf 1\{y < s\}\,\dd s$, the denominator with the soft cutoff is the average, over independent thresholds $s_b \in [1, 1 + \kappa_s]$ of the bonds, of the denominators with the sharp cutoffs $\abs{(1/i)\log V'(b)} < s_bM_0\varepsilon_k$, and at real configurations it lies between the sharp denominators at $s_b = 1$ and $s_b = 1 + \kappa_s$. In the variable $B'$ of \cite[(1.1), (1.2), pp.~356--357]{BalabanLFR2} the sharp cutoff is $\abs{B'} < s\,t$ with $t \coloneqq M_0g_k^{-1}\varepsilon_k$. The difference of the two denominators is an integral over the set where $t \leq \abs{B'(b)} < (1 + \kappa_s)t$ for some bond $b$; there the real part of the quadratic form is at least $c_\Lambda\norm{B'}^2$ with a constant $c_\Lambda > 0$, by \cite[(1.7)--(1.9), p.~358]{BalabanLFR2}, and the remaining factors of the integrand and the denominator at $s = 1$, bounded below by \cite[(1.10), (1.11), p.~358]{BalabanLFR2}, contribute at most a factor $e^{C(M)}$, since $\Lambda$ has at most $O(M^4)$ bonds. So the two denominators differ relatively by at most $O(M^4)e^{C(M)}\exp(-\frac12c_\Lambda t^2)$. Since $t^2 = M_0^2\varepsilon_k^2g_k^{-2} \geq \rho^{-1}M_0^2A_0^2(\log g_k^{-2})^{2p_0}$ by Paper, Lemma~\ref{P-lem:tube-couplings}(i), this is at most $\frac14$ for $\gamma$ small, and the logarithms of the denominators differ by at most $\frac12$.

(e) In the term without large-field regions every domain $\Omega_j$, $\Lambda_j$ is the whole torus, every test function of an earlier step has passed, and Ba{\l}aban removes each of them on the support of the test functions of the next step, \cite[(3.6)--(3.9), pp.~265--266]{BalabanCRE}, ``on $\Omega_{k+1}^{\sim}$ the functions $\chi_{Ax}$ and $\chi_k$ are equal to 1'' \cite[p.~265]{BalabanCRE}. By (a) these removals hold for the soft test functions, and by induction on the step only $\chi_k$ remains. No $\mathbf R$-operation acts on this term, since it has no large-field component.
\end{proof}

The proposition applies verbatim to the first step of the paper, \cite{WilsonFirstStep}, which is the step of \cite[Sect.~1]{BalabanCRE} with the loop inserted: its bounds drawn from the passed test functions (1.4) and (1.8) of \cite{BalabanCRE} hold at $(1 + \kappa_s)\varepsilon_1L^{-2}$ and $2(1 + \kappa_s)\delta_0$, and its uses of failed test functions, which exclude a label by a bound below the threshold and take a witness from a nonvanishing complement, hold as printed, since $1 - \bar\chi(f/t)$ is nonzero only for $f > t$.

\subsection{The Depth}\label{sec:depth}

Ba{\l}aban does not fix the depth $N$ of the $\mathbf R$-operation: ``conditions on $N$ will be formulated in constructions of this section'' \cite[p.~177]{BalabanLFR1}. In the large-field induction of \cite{BalabanLFR2} he takes $N = R_j$ \cite[p.~384]{BalabanLFR2}, and with this choice $L^N$ is not bounded by a power of $\log g_k^{-2}$, which the inductive step of the paper needs (Paper, Section~\ref{P-sec:depth}). The paper takes $N = N_k$ of Paper, Definition~\ref{P-def:depth}, at the $\mathbf R$-operation of the step $k$, and Paper, Lemma~\ref{P-lem:depth}, gives the bounds
\begin{equation}\label{eq:depth-recalled}
    L^{-N_k} \leq (\log g_k^{-2})^{-\nu_L}\;,
    \qquad
    L^{N_k} \leq C_{\mathrm{dep}}(\log g_k^{-2})^{\nu_L}\;,
    \qquad
    N_k \leq R_k\;,
    \qquad
    N_{k+1} \leq N_k\;,
\end{equation}
with $\nu_L = p_0 + q_1 + 7r + 1$, for every history of a tube. It remains to see that these are all that the two papers need.

\begin{proposition}[The Depth]\label{prop:depth}
Let $\gamma$ satisfy the hypotheses of Paper, Lemma~\ref{P-lem:depth}, and be small. For every tube and every history as in Proposition~\ref{prop:parameters}, every statement of \cite{BalabanLFR1,BalabanLFR2} in which the depth enters holds for the construction of the paper, in which the $\mathbf R$-operation of the step $k$ has the depth $N_k$. In the large-field induction \cite[(1.80), pp.~384--387]{BalabanLFR2}, stated for $N = R_j$, the number $K$ of steps which a domain needs before the $\mathbf R$-operation of the step $j + K$ removes it is at most its value for $N = R_j$, and the bound (1.80) holds as printed.
\end{proposition}

\begin{proof}
We go through the places at which $N$ enters, and we sort them into four kinds.

\emph{The conditions which Ba{\l}aban states.} They are (N1)--(N4) of Paper, Section~\ref{P-sec:depth}: $N > N_0$ \cite[p.~179]{BalabanLFR1}; the two restrictions of \cite[(1.94), p.~198]{BalabanLFR1}, $N \leq O(1)(\log g_k^{-2})^\nu$ with $\frac12\nu \leq p_0 - p_1 - 1$, and $N$ so large that the constant of the second term of (1.94) is at most $\frac12\alpha$; $N \leq R_k$ \cite[p.~361]{BalabanLFR2}; and $L^{-N} \leq (\log g_k^{-2})^{-\nu}$ for a sufficiently large $\nu$ \cite[p.~363]{BalabanLFR2}, which recurs as ``the conditions discussed above'' and ``the usual conditions on $N$'' \cite[p.~364, (1.47), p.~368]{BalabanLFR2} and as ``our assumptions on $N$'' \cite[p.~382]{BalabanLFR2}. Paper, Section~\ref{P-sec:depth}, shows that $N_k$ satisfies each of them, by \eqref{eq:depth-recalled}: in particular the product $O(1)A_0C_1B_3^3B_5M^{10}p_0(g_k)q_1(g_k)R_k^7L^{-N}$ of \cite[(1.47)]{BalabanLFR2} is at most a constant times $(\log g_k^{-2})^{p_0+q_1+7r-\nu_L} = (\log g_k^{-2})^{-1}$, and the number $O(1)(A_0/A_1)B_3^2B_5M^7R_j^{1+\beta_0}(p_0(g_j)/p_1(g_j))L^{-N}$ of \cite[p.~382]{BalabanLFR2}, which has to be below $\frac14$, is at most a constant times $(\log g_j^{-2})^{r(1+\beta_0)+p_0-p_1-\nu_L}$, whose exponent is negative. In the last bound the scale is $j$ and the depth that of the step $j$; at the $\mathbf R$-operation of a later step $k$ the depth $N_k \leq N_j$ is smaller, by \eqref{eq:depth-recalled}, and $L^{-N_k} \leq (\log g_k^{-2})^{-\nu_L} \leq (1 + \beta_0)(\log g_j^{-2})^{-\nu_L}$, by Ba{\l}aban's (2.7) for the history, so that the same bound holds with a constant larger by the factor $1 + \beta_0$.

\emph{Upper bounds.} The depth enters as an upper bound, through $N \leq R_k$, in $11d^2\varepsilon_h \leq 11d^2(1 + \beta_0)^2N^{\beta_0}\varepsilon_k \leq 11d^2(1 + \beta_0)^2R_k^{\beta_0}\varepsilon_k$ \cite[p.~362]{BalabanLFR2}, in $O(N) \leq O(R_k)$ and $N^{2+4\beta_0}(100MR_k)^4 \leq O(R_k^7)$ \cite[p.~366]{BalabanLFR2}, in $R_{j-N+1} \leq (L + 1)N^{\beta_0}R_j$ \cite[pp.~382, 383]{BalabanLFR2}, and in the range $j - N \leq m \leq j$ of the scales of a factor \cite[p.~381]{BalabanLFR2}. Each of these expressions increases with $N$, and each holds for $N_k \leq R_k$.

\emph{Step counts.} In the proof of (1.80) the number $K$ counts the steps after which the domain $S^K(Z)$ satisfies the conditions (i), (ii) of \cite[p.~177]{BalabanLFR1} with $N = R_j$: $K \leq n_0 - j + R_j$ \cite[p.~385]{BalabanLFR2}, $K = R_{j+1}$ for a small domain \cite[p.~386]{BalabanLFR2}, and $K \leq K_2 + n_1 + R_{j+1}$ \cite[p.~387]{BalabanLFR2}. The condition (ii) asks that no new large-field region has been created in the domain in the preceding $N$ steps, and it is weaker for a smaller $N$; the $\mathbf R$-operation which removes $S^K(Z)$ is that of the step $j + K$, with the depth $N_{j+K} \leq N_j \leq R_j$, by \eqref{eq:depth-recalled}. So the number of steps after which the domain is removed is at most the number $K$ of Ba{\l}aban's counts, and since the right side of (1.80) increases with $K$, the bound holds as printed.

\emph{The other uses.} In \cite{BalabanLFR1} the depth enters, besides (N1) and (1.94), the proof of the removal of $\chi''_k$: the bound $22d^2\max\{1, L^{-2N}N^{1/2}O(1)B_3B_5M^5(1 + \beta_0)\}\varepsilon_h + 2\delta'_k < 23d^2\varepsilon_h$ \cite[p.~199]{BalabanLFR1}, which needs $L^{-2N}N^{1/2}O(1)B_3B_5M^5(1 + \beta_0) \leq 1$ and $N^{1/2}p_1(g_h)/p_0(g_h)$ small, and the bound $B_3(1 + \beta_0)(1 + (k - j)^{1/2})(A_1/A_0)(p_1(g_j)/p_0(g_j)) < \alpha$ for $h \leq j \leq k$, where $k - j \leq N$ \cite[p.~199]{BalabanLFR1}. Since $p_1(g)/p_0(g) = A_0^{-1}(\log g^{-2})^{p_1-p_0}$ with $p_0 - p_1 > 1$, and $(\log g_j^{-2})^{p_1-p_0} \leq (1 + \beta_0)(\log g_k^{-2})^{p_1-p_0}$ for $j \leq k$ by Ba{\l}aban's (2.7), and since $N_k \leq \log_LC_{\mathrm{dep}} + \nu_L\log_L\log g_k^{-2}$ by \eqref{eq:depth-recalled}, both products are at most a constant times $(\log\log g_k^{-2})^{1/2}(\log g_k^{-2})^{p_1-p_0}$, which tends to zero; and $L^{-2N_k}N_k^{1/2} \to 0$ as well. The depth also enters definitions: the powers $L_0^{2N_0}$ of \cite[p.~192]{BalabanLFR1}, the width of a boundary layer \cite[p.~363]{BalabanLFR2}, the number of the new boundary terms \cite[p.~369]{BalabanLFR2}, and the spaces of the analytic extensions at $n = N$ \cite[p.~375]{BalabanLFR2}; these hold at every $N$.

Every place at which the depth enters the two papers is of one of these kinds, and in each case the statement holds for $N_k$.
\end{proof}

\section{The Coefficients of the Flow}\label{sec:coefficients}

Paper, Proposition~\ref{P-prop:coefficients}, asserts three properties of the coefficient $\beta_{k+1}(h)$ of the step $k + 1$ along a history $h$ of a tube: it differs from the coefficient $b_{k+1}$ at zero coupling by at most $\delta'$, it differs from it by a multiple of $g_k^2$, and it is continuous in $h$. By Paper, \eqref{P-eq:beta-split}, the difference is the polarisation coefficient $\Lambda$ of the fluctuation integral of the step, and two things stand between this and the bounds. Ba{\l}aban bounds the terms of the fluctuation integral by a multiple of the threshold of the fluctuation field \cite[(2.41), p.~21]{BalabanRG2}, which gives (i) once $\Lambda$ is shown to be bounded on local terms, but a bound by $g_k^2$ needs a symmetry which he does not use: replaced by a complex number $z$ in the exponent, the coupling enters only through $z$ times the field, and the integral is an even function of $z$. The threshold breaks this picture, since the domain of analyticity in $z$ shrinks as the threshold grows, and the second idea is to telescope over thresholds which double, the differences being Gaussian tails. The $\mathbf R$-terms of the earlier steps enter the fluctuation integral as well, and they are separated by a family of densities analytic in a factor multiplying them. This section proves (i) and (ii), as Proposition~\ref{prop:coefficient-bounds}; the continuity (iii) is Proposition~\ref{prop:coefficient-continuity}, and the coefficients at zero coupling are the subject of Proposition~\ref{prop:one-loop-window}.

\subsection{The Fluctuation Integral as a Function of the Coupling}\label{sec:fluctuation-integral}

Fix a tube, an integer $k < K$, and a history $h = (x_0, \dots, x_k) \in H_k$ with $g_k = x_k^{-1/2}$. We use the units of the step $k + 1$ and the notation of \cite[Sect.~2]{BalabanRG1} and \cite{BalabanRG2}. The fluctuation integral of the step is the logarithm of the integral \cite[(2.12), (2.13), p.~268]{BalabanRG1} of the exponent of \cite[(2.12)]{BalabanRG1} against the Gaussian measure $\dd\mu_{C^{(k)}}(B)$, with the cutoff of the fluctuation field. In the construction of the paper the cutoff restricts the unscaled field by the threshold $\vartheta_k$ (Paper, Definition~\ref{P-def:construction}(ii)); over the bonds $b \notin \{ b_0(c) \}$ on which \cite[(2.9), p.~266]{BalabanRG1} acts it depends on the Gaussian variable only through $\abs{B(b)} < p$, where
\begin{equation}\label{eq:scaled-threshold}
    p \coloneqq \frac{\vartheta_k}{g_k} = \frac{\hat g_k}{g_k}A_1(\log\hat g_k^{-2})^{p_0} \geq \rho^{-1/2}A_1(\log g_k^{-2})^{p_0}\;,
\end{equation}
by Paper, Definition~\ref{P-def:tube-parameters}(ii) and Lemma~\ref{P-lem:tube-couplings}(i). For a threshold $q > 0$ let $\chi_q$ be the product over the same bonds of the characteristic functions of $\{ \abs{B(b)} < q \}$. For $z \in \mathbb C$ let $W(z, B) \coloneqq \sum_YV_k(Y; z, B)$ be the localised exponent \cite[(1.41), p.~11]{BalabanRG2} with $g_k$ replaced by $z$, and put
\begin{equation}\label{eq:complex-fluctuation-integral}
    E(z; q) \coloneqq \log\int\chi_q\exp W(z, B)\,\dd\mu_{C^{(k)}}(B)\;,
\end{equation}
so that the fluctuation part of the regular term of the step is $\mathbf E^{(k+1)}_{\mathrm{fl}} = E(g_k; p)$.

\emph{Parity.} Every term of the exponent of \cite[(2.12), p.~268]{BalabanRG1} is a function of $g_kCB$, or $g_k^{-2}$ times such a function of order at least three: the terms $\operatorname{Tr}\log(I - h(\delta\tilde D/\delta B)(g_kCB))$ and $\log\sigma(g_kCB - h\tilde D(g_kCB))$, and the terms $g_k^{-2}\mathbf G_3(g_kCB)$ and $g_k^{-2}\langle H_1g_kCB, \Delta_1H_1h\tilde D(g_kCB)\rangle$, $\tilde D$ starting at order two. The localisation of \cite[(1.33)--(1.41), pp.~9--11]{BalabanRG2} keeps this form when $g_k$ is replaced by $z$: its pieces are functions of $zCB$, with the prefactor $z^{-2}$ or $1$, and the bracket of the effective action enters through $(z/g_k)B$ at the real coupling (Lemma~\ref{lem:complex-coupling}). Hence $W(-z, B) = W(z, -B)$. The Gaussian measure is centred and $\chi_q$ is even, so that
\begin{equation}\label{eq:parity}
    E(-z; q) = E(z; q)
\end{equation}
as functions of the background field and the current.

\begin{lemma}[The Polarisation Is Bounded]\label{lem:polarisation-bounded}
Let $F = \sum_XF(X)$, with $F(X)$ analytic on the space $U^c_{k+1}(X, \alpha_0, \alpha_1)$ of \cite[(1.10)--(1.16), p.~262]{BalabanRG1}, depending on the configuration restricted to $X$, and with $\abs{F(X)} \leq Ae^{-\kappa d_{k+1}(X)}$. Let $r \in (0, 1]$ be such that the polydisc of the configuration of radius $r$ in the norms of the space, about the trivial configuration, lies in it. Then the polarisation coefficient \eqref{P-eq:polarisation-coefficient} of the paper satisfies $\abs{\Lambda(F)} \leq C_\Lambda(1)r^{-2}A$, where $C_\Lambda(1)$ depends on $\kappa$, $\delta_0$, $d$, $M$, and $L$ only.
\end{lemma}

\begin{proof}
Write $u(B)$ for the configuration $(U_{k+1}(e^{iB}), J_{k+1}(e^{iB}))$ restricted to $X$. The map $B \mapsto U_{k+1}(e^{iB})$ is analytic for small complex $B$, and its derivative decays exponentially with the rate $\delta_0$, by \cite[Proposition~9, p.~309]{BalabanVAR} and \cite[(190), p.~308]{BalabanVAR}, so that $\norm{\partial u/\partial B_\nu(y)} \leq Ce^{-\delta_0\dist(y, X)}$; the bound (189) from which \cite[p.~308]{BalabanVAR} derives (190) is stated there without the calculation, and Lemma~\ref{lem:var-189} proves it. Cauchy's estimate in a single variable $B_{\nu'}(y')$ on a disc of fixed radius, applied to this derivative and symmetrised, gives
\begin{equation}\label{eq:second-derivative-background}
    \Bigl\lVert\frac{\partial^2u}{\partial B_\nu(y)\partial B_{\nu'}(y')}\Bigr\rVert \leq Ce^{-\delta_0(\dist(y, X) + \dist(y', X))/2}\;.
\end{equation}
Cauchy's estimate on the lines $u + tv$ and $u + tv + sw$ in the configuration gives $\abs{DF(X)[v]} \leq Ae^{-\kappa d_{k+1}(X)}\norm{v}/r$ and $\abs{D^2F(X)[v, w]} \leq 4Ae^{-\kappa d_{k+1}(X)}\norm{v}\norm{w}/r^2$. By the chain rule the second derivative $\Pi(x)$ of Paper, \eqref{P-eq:polarisation-coefficient}, is the sum over $X$ of $D^2F(X)[\partial_{B_\mu(0)}u, \partial_{B_\nu(x)}u] + DF(X)[\partial^2_{B_\mu(0)B_\nu(x)}u]$, and with $r \leq 1$ the term of $X$ is at most $CAr^{-2}e^{-\kappa d_{k+1}(X)}e^{-\delta_0(\dist(0, X) + \dist(x, X))/2}$. Finally $\sum_x\abs{x}^2e^{-\delta_0\dist(x, X)/2} \leq C(1 + \diam X + \dist(0, X))^{d+2}$, and $\diam X \leq M(2d_{k+1}(X) + 4)$, by Lemma~\ref{lem:cube-count}; summing first over the $X$ which contain a given cube, the cube of $X$ nearest to $0$, by \cite[(1.26), p.~8]{BalabanRG2}, and then over the cube, the series converges for $\kappa$ large, and its sum is linear in $A$ and proportional to $r^{-2}$.
\end{proof}

\begin{lemma}[Complex Coupling]\label{lem:complex-coupling}
Let $\varepsilon^*$ be such that $A_\chi\varepsilon^*$ satisfies every condition which \cite{BalabanRG1,BalabanRG2} place on their cutoff $\varepsilon_1$, where $A_\chi \coloneqq (1 + 2^{12}(M + 1)^4)(1 + 1/(\delta_{\mathrm{RG}}\kappa))$ and $\delta_{\mathrm{RG}} \coloneqq (1 - 2L^{-1})/12$, and let $q_{\mathrm{th}}$ be the number of Lemma~\ref{lem:threshold-differences}. Let $q \geq q_{\mathrm{th}}$ and $\abs{z}q < \varepsilon^*$. Then the expansion \cite[(2.1)--(2.13), pp.~12--14]{BalabanRG2} applies to \eqref{eq:complex-fluctuation-integral} and gives $E(z; q) = \sum_XE(X; z, q)$, where each $E(X; z, q)$ is analytic in $z$ on $\abs{z} < \varepsilon^*/q$ and in the configuration on $U^c_{k+1}(X, \alpha_0, \alpha_1)$, satisfies
\begin{equation}\label{eq:complex-step-bound}
    \abs{E(X; z, q)} \leq O(1)C_3\varepsilon^*\exp(-\kappa d_{k+1}(X))\;,
\end{equation}
with the constant $C_3$ of \cite[p.~20]{BalabanRG2}, and equals Ba{\l}aban's term at $z = g_k$, $q = p$.
\end{lemma}

The conditions on $\varepsilon_1$ are all upper bounds on it, so that $\varepsilon^*$ exists; they are listed in Lemma~\ref{lem:rg2-complex}. The number $\delta_{\mathrm{RG}}$ replaces Ba{\l}aban's $\delta = (1 - 2L^{-1})/10$ \cite[p.~21]{BalabanRG2}, which he uses only to bound $\exp(-(1 - 10\delta)\frac12L\kappa d_{k+1})$ by $\exp(-\kappa d_{k+1})$ \cite[(2.41), p.~21]{BalabanRG2}; with the smaller value, $(1 - 12\delta_{\mathrm{RG}})\frac12L = 1$ leaves room for one more loss of $\delta_{\mathrm{RG}}$ in the exponent, which Lemma~\ref{lem:threshold-differences} uses, and his other conditions on $\delta$ are lower bounds on $\kappa$.

\begin{proof}
The bracket of the effective action in the exponent of \cite[(2.12)]{BalabanRG1} depends on $(g_k, B)$ only through $g_kCB$, and \cite[(1.33)]{BalabanRG2} localises it at real $B$ as $\sum_YV'_k(Y, B)$. The bracket term of $W$ is $\sum_YV'_k(Y, (z/g_k)B)$, the same sum at the real coupling $g_k$ and a complex field; $V'_k$ is holomorphic on the polydisc $\abs{(z/g_k)B} < \varepsilon^*/g_k$ and satisfies \cite[(1.36), p.~9]{BalabanRG2} there, by Lemma~\ref{lem:rg2-complex}, so that $V'_k(Y, (z/g_k)B)$ is analytic in $z$ for $\abs{z}\abs{B} < \varepsilon^*$ on $Y$. The other terms of the exponent are built in \cite[(1.38)--(1.40), p.~10]{BalabanRG2} from the unscaled fluctuation field $B' = g_kCB - h\tilde D(g_kCB)$ \cite[(1.19), p.~6]{BalabanRG2}, and ``it can be done in the same way for all terms'' \cite[p.~11]{BalabanRG2}; with $z$ in place of $g_k$ each of them is the same analytic function of $zCB$ and of $B' = zCB - h\tilde D(zCB)$. The prefactor $z^{-2}$ cancels: the second-order expansion \cite[(1.40)]{BalabanRG2} writes the pieces as $z^{-2}\langle\hat QB', B'\rangle = \langle\hat QB'/z, B'/z\rangle$, and $B'/z = CB - z^{-1}h\tilde D(zCB)$ is analytic also at $z = 0$, since $\tilde D(zCB)$ is holomorphic in $z$ and of order $\abs{z}^2\abs{CB}^2$; the other terms with the prefactor $z^{-2}$ are of order at least three in $zCB$. By \cite[(1.20), p.~6]{BalabanRG2} with $\abs{z}$ in place of $g_k$, $\abs{B'} \leq C'_1\abs{z}\abs{B} < C'_1\varepsilon^*$, so that \cite[(1.43), (1.36)]{BalabanRG2} hold with $\varepsilon^*$ in place of $\varepsilon_1$ and each $V_k(Y; z, B)$ is analytic in $z$ for $\abs{z}\abs{B} < \varepsilon^*$ on $Y$; Lemma~\ref{lem:rg2-complex} carries Ba{\l}aban's Lemmata 1 and 2, stated for real $g_k$, to this setting.

The expansion of \cite[Sect.~2]{BalabanRG2} uses the potentials only through \cite[(1.42), (1.43), (1.36)]{BalabanRG2} and the count of the cubes, in \cite[(2.18)--(2.20), p.~16]{BalabanRG2}, and through their analyticity in the configuration \cite[(1.34), p.~9]{BalabanRG2}, which does not involve $z$. It uses the cutoff only through a factor $\chi_{k,Y_0}\chi^c_{k,P}$ with values in $[0, 1]$ and the tail bound \cite[(2.22), p.~16]{BalabanRG2}: the potentials are evaluated on the support of $\chi_{k,Y_0}$ \cite[(2.3), p.~12]{BalabanRG2}, where $\abs{z}\abs{B} < \abs{z}q < \varepsilon^*$, and the threshold enters the conditions of the expansion only as $\varepsilon_1/g_k$ in \cite[(2.22), (2.31), (2.33), (2.34)]{BalabanRG2}, which with $q$ in its place are lower bounds on $q$ and hold for $q \geq q_{\mathrm{th}}$. So \cite[(2.1)--(2.41)]{BalabanRG2} hold with $\varepsilon^*$ in place of $\varepsilon_1$ and $q$ in place of $\varepsilon_1/g_k$, which gives \eqref{eq:complex-step-bound}. Each term \cite[(2.14), p.~15]{BalabanRG2} is an integral of a function analytic in $z$, dominated by \cite[(2.26), p.~17]{BalabanRG2} uniformly on compact subsets of $\abs{z} < \varepsilon^*/q$, and the sums \cite[(2.35)--(2.41)]{BalabanRG2} converge uniformly there, so that $E(X; z, q)$ is analytic in $z$.
\end{proof}

\begin{lemma}[Threshold Differences]\label{lem:threshold-differences}
There is a number $q_{\mathrm{th}}$, depending on $\varepsilon^*$, $\kappa$, $M$, and the constants of \cite{BalabanRG2}, such that the following holds. Let $q_{\mathrm{th}} \leq q' < q'' \leq 2q'$ and $\abs{z}q'' < \varepsilon^*$. Then $E(z; q'') - E(z; q') = \sum_X\Delta(X; z)$, where each $\Delta(X; z)$ is analytic in $z$ and in the configuration on $U^c_{k+1}(X, \alpha_0, \alpha_1)$, and
\begin{equation}\label{eq:threshold-difference}
    \abs{\Delta(X; z)} \leq C_4\exp(-\tfrac14\gamma_2q'^2)\exp(-\kappa d_{k+1}(X))\;,
    \qquad
    \abs{E(X; 0, q)} \leq C_5\exp(-\tfrac{1}{20}\gamma_2q^2)\exp(-\kappa d_{k+1}(X))
\end{equation}
for $q \geq q_{\mathrm{th}}$, with constants $C_4$, $C_5$ which do not depend on $k$, the history, $q'$, $q''$, $q$, or $z$, and with the constant $\gamma_2$ of \cite[(2.22)]{BalabanRG2}.
\end{lemma}

\begin{proof}
The families $\mathbf D$, the sets $P \subset Y_0^{c*}$, and the domains $Z_0$ and $Z$ of \cite[(2.1)--(2.13)]{BalabanRG2} do not depend on the threshold \cite[(2.1)--(2.3), p.~12]{BalabanRG2}, so that the activities at $q''$ and at $q'$ differ term by term only in the factor $\chi = \chi_{Y_0}\chi^c_P$. With $\chi''$ and $\chi'$ this factor at the two thresholds,
\begin{equation}\label{eq:cutoff-difference}
    \chi'' - \chi' = (\chi''_{Y_0} - \chi'_{Y_0})\chi''^c_P + \chi'_{Y_0}(\chi''^c_P - \chi'^c_P)\;.
\end{equation}
The first term is nonzero only if some bond $b$ of $\chi_{Y_0}$ has $\abs{B(b)} \geq q'$, and the second only if $P \neq \emptyset$ and every $b \in P$ has $\abs{B(b)} \geq q'$. Bounding each indicator as in \cite[(2.22)]{BalabanRG2} by $\mathbf 1\{ \abs{B(b)} \geq q' \} \leq \exp(-\frac12\gamma_2q'^2 + \frac12\gamma_2\abs{B(b)}^2)$, the first term by a sum over its bonds, and using $\frac12(1 + \abs{P}) \geq \frac14 + \frac14\abs{P}$ and $\frac12\abs{P} \geq \frac14 + \frac14\abs{P}$ for $P \neq \emptyset$, we obtain
\begin{equation}\label{eq:cutoff-difference-bound}
    \abs{\chi'' - \chi'} \leq (N_\chi(Y_0) + 1)\exp(-\tfrac14\gamma_2q'^2)\exp(-\tfrac14\gamma_2q'^2\abs{P} + \tfrac12\gamma_2\norm{Z_0B}^2)\;,
\end{equation}
where $N_\chi(Y_0)$ is the number of the bonds of $\chi_{Y_0}$. The factor $\exp(\frac12\gamma_2\norm{Z_0B}^2)$ is absorbed into the Gaussian measure as in \cite[(2.23), p.~16]{BalabanRG2}, and the factor $\exp(-\frac14\gamma_2q'^2\abs{P})$ runs the sums \cite[(2.31)--(2.34)]{BalabanRG2} with $\gamma_2/2$ in place of $\gamma_2$, which hold for $q' \geq q_{\mathrm{th}}$ once $q_{\mathrm{th}}$ is large. Each bond of $\chi_{Y_0}$ has an endpoint in $Y_0 = \bigcup_{Y \in \mathbf D}Y$, a closed cube of side $M$ contains at most $(M + 1)^4$ lattice points, each of which is an endpoint of $8$ bonds, and $M^{-4}\abs{Y} \leq 2^9(1 + d_k(Y))$ by Lemma~\ref{lem:cube-count}, so that
\begin{equation}\label{eq:bond-count}
    N_\chi(Y_0) + 1 \leq 1 + \sum_{Y \in \mathbf D}2^{12}(M + 1)^4(1 + d_k(Y)) \leq \prod_{Y \in \mathbf D}A_\chi e^{\delta_{\mathrm{RG}}\kappa d_k(Y)}\;,
\end{equation}
by $1 + \sum_Ya_Y \leq \prod_Y(1 + a_Y)$ and $1 + s \leq (1 + 1/(\delta_{\mathrm{RG}}\kappa))e^{\delta_{\mathrm{RG}}\kappa s}$ for $s \geq 0$. Each factor joins the factor of $Y$ in \cite[(2.26), p.~17]{BalabanRG2}, which becomes $2E_0A_\chi\varepsilon^*C_1\alpha_4^{-1}M^qe^{C_2\kappa_1}\exp(-(1 - 4\delta_{\mathrm{RG}})\kappa d_k(Y))$; since $A_\chi\varepsilon^*$ satisfies the condition of \cite[p.~18]{BalabanRG2}, the chain \cite[(2.27)--(2.38)]{BalabanRG2} runs with $A_\chi\varepsilon^*$ in place of $\varepsilon_1$ and one more $\delta_{\mathrm{RG}}$ in every exponent, and bounds the difference of the activities by its bound times $O(1)\exp(-\frac14\gamma_2q'^2)$. In the polymer series \cite[(2.13)]{BalabanRG2} the difference of two products is a sum of products with one factor of the difference, and \cite[(2.39)--(2.41)]{BalabanRG2} give the first bound of \eqref{eq:threshold-difference}, the extra $\delta_{\mathrm{RG}}$ being absorbed by $(1 - 12\delta_{\mathrm{RG}})\frac12L = 1$.

At $z = 0$ every localised potential vanishes: in Ba{\l}aban's Lemma~1 the function $H_k(s, 0) = 0$ \cite[(1.18), p.~6]{BalabanRG2}, so that the Cauchy integral which defines the potential vanishes; in the action terms $B' = 0$ and the coefficients of \cite[(1.40)]{BalabanRG2} are of order $\abs{B'}$; and the terms without the prefactor, $\operatorname{Tr}\log(I - h(\delta\tilde D/\delta B))$ and $\log\sigma$, vanish at $B' = 0$, because $\tilde D$ starts at order two and $\sigma$ there is normalised to one at $B' = 0$. So only $\mathbf D = \emptyset$ contributes to \cite[(2.1)]{BalabanRG2}, every nonzero term has $P \neq \emptyset$, and \cite[(2.33)]{BalabanRG2} gives the factor $\exp(-\frac{1}{20}\gamma_2q^2)$ per component in place of $\max\{ \varepsilon_2, \dots \}$, which with \cite[(2.37)--(2.41)]{BalabanRG2} is the second bound of \eqref{eq:threshold-difference}.
\end{proof}

\begin{lemma}[The $\mathbf R$-Terms]\label{lem:r-term-shift}
Let $\gamma$ be small, let $\beta_{k+1}(h)$ be the coefficient of the construction of the paper along $h$, and let $\beta^0_{k+1}(h)$ be the coefficient of the same recursion along $h$ in which no $\mathbf R$-term enters the regular part. Then
\begin{equation}\label{eq:r-term-shift}
    \abs{\beta_{k+1}(h) - \beta^0_{k+1}(h)} \leq C_Rg_k^{\kappa_0-4}\;,
\end{equation}
with a constant $C_R$ which does not depend on $k$, $K$, the history, or the tube, and with the exponent $\kappa_0$ of \cite[(2.31), p.~260]{BalabanCRE}, which Ba{\l}aban takes $\geq 7$ \cite[(2.46), p.~263]{BalabanCRE}.
\end{lemma}

\begin{proof}
The $\mathbf R$-terms satisfy $\abs{\mathbf R^{(j)}(X, (\mathbf U, \mathbf J))} \leq g_j^{\kappa_0}\exp(-\kappa d_j(X))$ by Proposition~\ref{prop:parameters}(e), and Ba{\l}aban includes their part $R'_k$ into the regular term of the step \cite[p.~271]{BalabanCRE}, while the part which stays an $\mathbf R$-term needs no counterterm \cite[pp.~260, 276]{BalabanCRE}. For $s \in \mathbb C$ we run the regular recursion of the construction along $h$ with $sR'$ in place of $R'$ at every step, and with the coefficients $\beta_j(s)$ computed from its own regular terms by \cite[(3.57), (3.61), pp.~281--282]{BalabanCRE}. At $s = 1$ this is the construction, and at $s = 0$ no $\mathbf R$-term enters, so that $\beta_{k+1}(1) = \beta_{k+1}(h)$ and $\beta_{k+1}(0) = \beta^0_{k+1}(h)$.

\emph{The size of the $\mathbf R$-part.} By \eqref{P-eq:tube-window} of the paper, $x_j \geq x_k + b(k - j) - C_W \geq \frac12x_k + (k - j)c\ln L$ for $j \leq k$, since $b \geq c\ln L$ and $C_W \leq t/8 \leq x_k/7$. So
\begin{equation}\label{eq:r-sum}
    \sum_{j \leq k}g_j^{\kappa_0} \leq \sum_{m \geq 0}\bigl(\tfrac12x_k + mc\ln L\bigr)^{-\kappa_0/2} \leq C_\Sigma g_k^{\kappa_0-2}\;,
    \qquad
    C_\Sigma \coloneqq 2^{\kappa_0/2}\Bigl(1 + \frac{2}{(\kappa_0 - 2)c\ln L}\Bigr)\;.
\end{equation}
The $\mathbf R$-part of the input of a step enters the fluctuation integral through the potential $V'$ of Ba{\l}aban's Lemma~1, which is built linearly from the bracket of the effective action \cite[(1.33)]{BalabanRG2}, the sums over the domains and the scales being controlled by the exponential factor and by $g_j^{\kappa_0}$ \cite[p.~260]{BalabanCRE}. So there is a constant $K_V$, depending on the constants of the step and on $\kappa$ only, such that an input whose $\mathbf R$-terms are bounded with $a_ig_i^{\kappa_0}$ in place of $g_i^{\kappa_0}$ enters with $E_0 + K_V\sum_{i \leq j}a_ig_i^{\kappa_0}$ in place of $E_0$ in \cite[(1.36), p.~9]{BalabanRG2}. Here $a_i = \abs{s}$. With $\theta \coloneqq E_0/(K_VC_\Sigma)$ and $\Theta_k \coloneqq \theta g_k^{2-\kappa_0}$, \eqref{eq:r-sum} gives $K_V\abs{s}\sum_{i \leq j}g_i^{\kappa_0} \leq E_0$ for $\abs{s} \leq \Theta_k$ and every $j \leq k$, so that the input of every step is bounded with $2E_0$ in place of $E_0$.

\emph{The step with $2E_0$ and complex $s$.} The step accepts such inputs. Its conditions which involve $E_0$ are upper bounds on the cutoff relative to the radii of the domains, \cite[pp.~18, 20, 21]{BalabanRG2}, which are imposed for $2E_0$ when the constants of the radii are chosen, and two lower bounds on the threshold which only get weaker as $E_0$ grows \cite[pp.~16, 19]{BalabanRG2}. The expansion of the step \cite[Sect.~2]{BalabanRG2}, which \cite[p.~277]{BalabanCRE} uses, is analytic in complex multipliers of the interaction: \cite[(2.14), p.~15]{BalabanRG2} carries complex coefficients $\tau(Y)$ and is bounded through the modulus of the interaction, the only positivity used being that of the Gaussian measure at the real configuration with zero current \cite[p.~15]{BalabanRG2}. A complex factor $s$ on the $\mathbf R$-part enters in the same way, so that by induction over the steps $j \leq k$ every regular term of the family satisfies \cite[(1.18), p.~263]{BalabanRG1} with $E_0$ and is analytic in $s$ on $\abs{s} < \Theta_k$.

\emph{Polarisation and Cauchy's estimate.} The fluctuation term of the step $k + 1$ is analytic on $U^c_{k+1}(X, \alpha_{0,k+1}, \alpha_{1,k+1})$, with the radii of the tube, and bounded by $E_0e^{-\kappa d_{k+1}(X)}$. The polydisc about the trivial configuration of radius $r \coloneqq \min(\alpha_{0,k+1}, \alpha_{1,k+1})/C_{\mathrm{cen}}$ lies in this space, with a constant $C_{\mathrm{cen}}$ depending on $L$, $M$, and the averaging operations only, and $r \leq 1$ for $\gamma$ small. Lemma~\ref{lem:polarisation-bounded} gives, with $\hat x \coloneqq \hat g_{k+1}^{-2} \leq \hat g_k^{-2} \leq \rho x_k \leq 2x_k$ and $q_* \coloneqq \min(q_0, q_1) \geq 2$,
\begin{equation}\label{eq:polarisation-tube}
    \abs{\beta_{k+1}(s) - b_{k+1}} \leq C_\Lambda(1)r^{-2}E_0 \leq C'E_0\hat x(\log\hat x)^{-2q_*} \leq 2C'E_0g_k^{-2}\;,
\end{equation}
with $C' \coloneqq C_\Lambda(1)C_{\mathrm{cen}}^2/\min(C^\alpha_0, C^\alpha_1)^2$, and $s \mapsto \beta_{k+1}(s)$ is analytic on $\abs{s} < \Theta_k$, the infinite-volume limit of $\Lambda$ being uniform in $s$ by the domination in Lemma~\ref{lem:polarisation-bounded}. For $\gamma$ small, $\Theta_k \geq \theta\gamma^{2-\kappa_0} > 2$, because $\kappa_0 > 2$ and $g_k \leq \gamma$, so that for $s_0 \in [0, 1]$ the closed disc of radius $\Theta_k/2$ about $s_0$ lies in $\abs{s} < \Theta_k$, and Cauchy's estimate gives $\abs{\beta'_{k+1}(s_0)} \leq 4C'E_0g_k^{-2}/\Theta_k$. Integrating over $[0, 1]$,
\begin{equation}\label{eq:r-shift-final}
    \abs{\beta_{k+1}(1) - \beta_{k+1}(0)} \leq \frac{4C'E_0g_k^{-2}}{\Theta_k} = 4C'K_VC_\Sigma g_k^{\kappa_0-4} \eqqcolon C_Rg_k^{\kappa_0-4}\;.
\end{equation}
The constants $C_\Sigma$, $K_V$, and $C'$ do not depend on the tube, so neither does $C_R$.
\end{proof}

\begin{proposition}[Bounds on the Coefficients]\label{prop:coefficient-bounds}
There are constants $C_\Lambda$, $C_3$, $C_{\mathrm{rem}}$, $C_R$, and $\gamma_1 > 0$, which do not depend on the tube, on $k$, or on the history, such that for every $\gamma \in (0, \gamma_1]$, every tube, every $k < K$, and every $h \in H_k$,
\begin{equation}\label{eq:coefficient-bounds}
    \abs{\beta_{k+1}(h) - b_{k+1}} \leq C_\Lambda C_3\gamma p_0(\gamma) + C_R\gamma^{\kappa_0-4}\;,
    \qquad
    \abs{\beta_{k+1}(h) - b_{k+1}} \leq \bigl(C_{\mathrm{rem}} + C_R\gamma^{\kappa_0-6}\bigr)g_k^2\;.
\end{equation}
\end{proposition}

\begin{proof}
By Lemma~\ref{lem:r-term-shift} it suffices to prove the bounds $C_\Lambda C_3\gamma p_0(\gamma)$ and $C_{\mathrm{rem}}g_k^2$ for $\beta^0_{k+1}(h)$, since $C_Rg_k^{\kappa_0-4} \leq C_R\gamma^{\kappa_0-4}$ and $C_Rg_k^{\kappa_0-4} \leq C_R\gamma^{\kappa_0-6}g_k^2$, as $g_k \leq \gamma$ by Paper, Lemma~\ref{P-lem:tube-couplings}(i). By Paper, \eqref{P-eq:beta-split}, and the linearity of $\Lambda$, $\beta^0_{k+1}(h) - b_{k+1} = \Lambda(\mathbf E^{(k+1)}_{\mathrm{fl}})$, the part of the regular term of order zero being $\log Z^{(k)}$, whose coefficient is $b_{k+1}$. We write $R(z; q) \coloneqq \Lambda(E(\cdot; z, q))$, so that $\beta^0_{k+1}(h) - b_{k+1} = R(g_k; p)$.

\emph{The first bound.} The cutoff of the step is that of \cite[Sect.~2]{BalabanRG2} with $\vartheta_k$ in place of $\varepsilon_1$, since $p = \vartheta_k/g_k$. The conditions which \cite{BalabanRG2} places on $\varepsilon_1$ are upper bounds, which hold for $\vartheta_k \leq \gamma p_0(\gamma)$ when $\gamma$ is small, and lower bounds on $\varepsilon_1/g_k$, which hold for $p$ by \eqref{eq:scaled-threshold}. So \cite[(2.41), p.~21]{BalabanRG2} bounds the terms of $\mathbf E^{(k+1)}_{\mathrm{fl}}$ by $O(1)C_3\vartheta_k\exp(-\kappa d_{k+1}(X))$; this is the ``additional factor $O(\varepsilon_k)$ arising from the bound of the fluctuation field'' of \cite[p.~277]{BalabanCRE}. Lemma~\ref{lem:polarisation-bounded}, with the radius of the polydisc fixed by the constants $\alpha_0$, $\alpha_1$ of \cite[(1.10)--(1.16), p.~262]{BalabanRG1}, gives $\abs{R(g_k; p)} \leq C_\Lambda O(1)C_3\vartheta_k$, where $C_\Lambda \coloneqq C_\Lambda(1)r^{-2}$, and we include the $O(1)$ in $C_3$. Finally $\vartheta_k = (A_1/A_0)\hat g_kp_0(\hat g_k) \leq \gamma p_0(\gamma)$, because $A_1 < A_0$, $\hat g_k \leq \gamma$, and $g \mapsto gp_0(g)$ increases on $(0, \gamma]$, its logarithmic derivative being $g^{-1}(1 - 2p_0/\log g^{-2}) > 0$ once $\log\gamma^{-2} > 2p_0$.

\emph{The second bound.} The function $R(\cdot; q)$ is analytic on $\abs{z} < \varepsilon^*/q$ and bounded there by $M_0 \coloneqq C_\Lambda O(1)C_3\varepsilon^*$, by Lemmata~\ref{lem:complex-coupling} and~\ref{lem:polarisation-bounded}, and it is even, by \eqref{eq:parity} and the linearity of $\Lambda$. An even function $f$ analytic on $\abs{z} < r_f$ with $\abs{f} \leq M_f$ has $(f(z) - f(0))/z^2$ analytic, and the maximum principle gives $\abs{f(g) - f(0)} \leq 2M_fg^2/r_f^2$ for $0 \leq g < r_f$. Let $J$ be the least integer with $2^Jq_{\mathrm{th}} \geq p$, put $p_j \coloneqq 2^jq_{\mathrm{th}}$ for $j < J$ and $p_J \coloneqq p$, and $D_j(z) \coloneqq R(z; p_{j+1}) - R(z; p_j)$. Then
\begin{equation}\label{eq:telescoping}
    R(g_k; p) = R(0; p) + \bigl[R(g_k; q_{\mathrm{th}}) - R(0; q_{\mathrm{th}})\bigr] + \sum_{j<J}\bigl[D_j(g_k) - D_j(0)\bigr]\;.
\end{equation}
The middle term is at most $2M_0q_{\mathrm{th}}^2g_k^2/\varepsilon^{*2}$. Each $D_j$ is even and analytic on $\abs{z} < \varepsilon^*/p_{j+1}$, which contains $g_k$ since $g_kp_{j+1} \leq g_kp = \vartheta_k < \varepsilon^*$, and bounded by $C_\Lambda C_4\exp(-\frac14\gamma_2p_j^2)$, by Lemmata~\ref{lem:polarisation-bounded} and~\ref{lem:threshold-differences}; with $p_{j+1} \leq 2p_j$ this gives $\abs{D_j(g_k) - D_j(0)} \leq 8C_\Lambda C_4p_j^2\exp(-\frac14\gamma_2p_j^2)g_k^2/\varepsilon^{*2}$, and the sum over $j$ is at most $8C_\Lambda C_4Sg_k^2/\varepsilon^{*2}$ with $S \coloneqq \sum_{j \geq 0}4^jq_{\mathrm{th}}^2\exp(-\frac14\gamma_24^jq_{\mathrm{th}}^2) < \infty$. The first term is $\Lambda$ of the logarithm of the Gaussian measure of the cutoff, and Lemma~\ref{lem:threshold-differences} at $z = 0$ gives $\abs{R(0; p)} \leq C_\Lambda C_5\exp(-\frac{1}{20}\gamma_2p^2) \leq C_\Lambda C_5g_k^2$, because $\frac{1}{20}\gamma_2p^2 \geq \frac{1}{20}\gamma_2\rho^{-1}A_1^2(\log g_k^{-2})^{2p_0} \geq \log g_k^{-2}$ for $g_k \leq \gamma$ small, by \eqref{eq:scaled-threshold}. So the second bound holds with
\begin{equation}\label{eq:c-rem}
    C_{\mathrm{rem}} \coloneqq \frac{2M_0q_{\mathrm{th}}^2}{\varepsilon^{*2}} + \frac{8C_\Lambda C_4S}{\varepsilon^{*2}} + C_\Lambda C_5\;.
\end{equation}
The constants are chosen in the order $\varepsilon^*$, $q_{\mathrm{th}}$, $\gamma_1$, and $\gamma_1$ is so small that for $g \leq \gamma_1$ the number $p$ of \eqref{eq:scaled-threshold} is at least $q_{\mathrm{th}}$, $gp < \varepsilon^*$, and the last inequality above holds. None of them depends on $k$, the history, or the tube.
\end{proof}

\subsection{Continuity in the History}\label{sec:continuity}

The coefficient $\beta_{k+1}(h)$ is the end product of $k + 1$ steps, each of which consists of integrals, quotients, cluster and localisation expansions, and the $\mathbf R$-operation, and the history enters all of them. Continuity in $h$ is therefore not a matter of inspecting a formula. Two features of the construction make it hold. The first is that no set on which a characteristic function, a test function, or a bump function is supported moves with the history, so that $h$ enters only through densities which are continuous in it and through the terms of the earlier steps. The second is analyticity: the localisation expansions read the earlier terms at complex arguments and through their derivatives, and the uniform bounds of the construction turn convergence at the real points into locally uniform convergence of all derivatives, by the theorems of Montel and Vitali. The coefficient is, finally, a limit as the torus grows, and we need this limit to be uniform in $h$, which Ba{\l}aban does not state; it comes from the random walk expansions which give the limit its existence.

One choice in the changes of variables has to be fixed for the first feature to hold. Ba{\l}aban specifies the bump function $g$ of \cite[(1.22), p.~251]{BalabanCRE} only on two sets \cite[p.~252]{BalabanCRE}, and a choice which depends on the coupling in any other way would move with the history.

\begin{definition}[The Bump Functions]\label{def:bump-functions}
Let $\varphi_1$ be a fixed $C^\infty$ function on $[0, \infty)$ with $0 \leq \varphi_1 \leq 1$, $\varphi_1 = 0$ on $[0, 4/3]$, and $\varphi_1 = 1$ on $[5/3, \infty)$. In the change of variables \cite[(1.22), p.~251]{BalabanCRE} of the step $1$, in its repetition at the step $k + 1$ \cite[p.~269]{BalabanCRE}, and in the change of variables \cite[p.~188]{BalabanLFR1} of the $\mathbf R$-operation, the bump function is
\begin{equation}\label{eq:bump-function}
    g(A') \coloneqq \varphi_1\Bigl(\frac{\abs{g_jA'}}{\vartheta}\Bigr)\;,
\end{equation}
where $g_j$ is the coupling of the scale of the change, $\vartheta$ its threshold ($\vartheta_0$, $\vartheta_k$, and $\delta'_j$ respectively), and $\abs{\cdot}$ the Hilbert--Schmidt norm \cite[p.~390]{BalabanPBG}.
\end{definition}

This choice has the properties which \cite[p.~252]{BalabanCRE} requires. It is $C^\infty$ on the Lie algebra, because $\varphi_1$ vanishes near $0$. It vanishes where $\abs{g_jA'} \leq \frac43\vartheta$. It equals one where $\abs{\exp ig_jA' - 1} \geq \frac53\vartheta$, because $\abs{\exp iB' - 1} \leq \abs{B'}$ for $B'$ in the Lie algebra: both matrices are diagonal in one orthonormal basis, with the entries $e^{i\lambda} - 1$ and $\lambda$, $\lambda$ real, and $\abs{e^{i\lambda} - 1} \leq \abs{\lambda}$. In the unscaled field $B' = g_jA'$ its support is fixed by the target.

\begin{proposition}[Continuity of the Coefficients]\label{prop:coefficient-continuity}
Let $\gamma \in (0, \gamma_1]$, with $\gamma_1$ as in Proposition~\ref{prop:coefficient-bounds}. For every tube and every $k < K$ the function $h \mapsto \beta_{k+1}(h)$ is continuous on $H_k$.
\end{proposition}

\begin{proof}
\emph{The claim of the induction.} A term $F(X)$ which the step $i$ produces is an analytic function of the configuration $(\mathbf U, \mathbf J)$ of the scale $i + 1$ on a domain $D_F$: the space $U^c_{i+1}(X, \alpha_{0,i+1}, \alpha_{1,i+1})$ for the regular terms \cite[(2.26)--(2.28), p.~259]{BalabanCRE} and for the $\mathbf R$-terms \cite[(2.31), p.~260]{BalabanCRE}, and the space \cite[(3.43), p.~276]{BalabanCRE}, which contains a fluctuation field $A$ with $\abs{A} < C_1p_1(\hat g_i)$, for the boundary terms. A \emph{real point} of $D_F$ is a point at which $\mathbf U$ is $G$-valued and $\mathbf J$ is real, and $A$ is real for a boundary term. The radii of these domains are those of the tube, fixed by the target, except the bound of the boundary terms on the unscaled field $g_iA$, which is $g_iC_1p_1(\hat g_i)$ and moves continuously with $h$; a real point at which this bound is not attained at $h^0$ is a real point for all $h$ near $h^0$. For $B$ real put $u_i(B) \coloneqq (U_{i+1}(e^{iB}), J_{i+1}(e^{iB}))$, call $B^0$ \emph{admissible} for $F$ if $u_i(B^0)$ is a real point of $U^c_{i+1}(X, \frac12\alpha_{0,i+1}, \frac12\alpha_{1,i+1})$, and let $P_i(B^0)$ be the complex polydisc $\{ \abs{B - B^0} < c_{\mathrm{pd}}\min(\alpha_{0,i+1}, \alpha_{1,i+1}) \}$. By \cite[(1.15), (1.16), p.~262]{BalabanRG1} and the bound on the derivative of the background field in the proof of Lemma~\ref{lem:polarisation-bounded}, there is $c_{\mathrm{pd}} > 0$ such that $u_i$ maps every $P_i(B^0)$ into $D_F$. The claim $S_j$, for $-1 \leq j \leq k$, is that the coefficients $\beta_i$, $i \leq j + 1$, are continuous on $H_j$, and that every term $F$ produced by the steps $0, \dots, j$ has the following two properties at every $h^0 \in H_j$:
\begin{enumerate}[label=(\alph*)]
\item $F(X, \mathbf U, \mathbf J; h) \to F(X, \mathbf U, \mathbf J; h^0)$ as $h \to h^0$, at every real point of $D_F$;
\item if $F$ is a regular term or an $\mathbf R$-term, then $F(X, u_i(B); h) \to F(X, u_i(B); h^0)$ uniformly on the compact subsets of $P_i(B^0)$, for every admissible $B^0$.
\end{enumerate}
A term depends on $h$ only through the entries of the history up to the step which produces it, and the later steps read it without changing it, so that (a) and (b), once proved at that step, hold at every later stage. The density of the step $0$ is $\exp[-x_0A - E]$, continuous in $x_0$ at fixed $U$, with no earlier terms, which is $S_{-1}$. We assume $S_{j-1}$ and prove $S_j$ in five steps.

\emph{Step 1: what depends on the history.} In the construction of the paper every threshold, cube size, depth, and partition, and every radius $\alpha_{i,j}$, is fixed by the target (Paper, Definition~\ref{P-def:construction}(i)). Every characteristic function is a function of the unscaled fields with a threshold fixed by the target: the functions $\chi_k(\Omega_k)$ of \cite[(2.17), p.~257]{BalabanCRE} on the plaquettes of the background fields \cite[(2.16)]{BalabanCRE}, which are defined from the fields alone; the function $\chi_{k,\Lambda}$ of \cite[(1.75), p.~193]{BalabanLFR1}; and the functions of the fluctuation fields, by Paper, Definition~\ref{P-def:construction}(ii), for example $\abs{\exp ig_jB_j(b) - 1} < 2\delta'_j$ in \cite[p.~188]{BalabanLFR1}. The soft test functions $\bar\chi(f/t)$ of Paper, Definition~\ref{P-def:construction}(iii), and those of Definition~\ref{def:soft-functions}, have $f$ a function of the unscaled fields and $t$ fixed by the target. The changes of variables \cite[(1.22)--(1.23), pp.~251--252]{BalabanCRE}, their repetition at \cite[p.~269]{BalabanCRE}, and that of \cite[p.~188]{BalabanLFR1} read, in the unscaled fields, $\exp iB = \exp iB'\exp(-ig(B')H)$, with the bump functions of Definition~\ref{def:bump-functions} and with $H$ an analytic function of the background fields which contains no coupling ($\mathbf H_{1,\mathrm{Ax}}$ of \cite[(1.21), p.~251]{BalabanCRE}, $\mathbf H^{(k)}$ of \cite[(3.22), p.~269]{BalabanCRE}, and $\mathbb H^{(j)}_Z$ of \cite[(1.58), p.~188]{BalabanLFR1}). These maps, their Jacobians \cite[(1.24), p.~252]{BalabanCRE}, the constraints of the block averages, and the Haar measures do not depend on $h$. The history therefore enters only through the Wilson action with the coefficient function \eqref{P-eq:history-coupling} of the paper, the Gaussian densities of the unscaled fluctuation fields, whose covariances are $g_j^2$ times operators which do not depend on $h$, the weight $\exp[-g_k^{-2}G]$ of the axial gauge and its normalisation, the Jacobians of the scalings, the densities and normalisations of the $\mathbf R$-operation \cite[(1.100), p.~201]{BalabanLFR1}, the counterterms $\beta_i$, and the terms of the earlier steps.

\emph{Step 2: where the step reads the earlier terms.} The terms of the step $j$ are defined by Ba{\l}aban's formulae as functions of the variables $(\mathbf U, \mathbf J)$ which replace $U_{j+1}$ and $J_{j+1}$ \cite[p.~272]{BalabanRG1}, and they read the earlier terms at three kinds of arguments:
\begin{enumerate}[label=(\roman*)]
\item in \cite[(3.13), p.~272]{BalabanRG1}, at $\exp i\xi\mathbf A\mathbf U$ and $(L^{j-1}\eta)^3\mathbf J + \Delta^\xi_{\mathbf U}\mathbf A - P_1\mathbf A + \mathbf F_1(\mathbf U, \mathbf A)$, for $\mathbf A$ in the class \cite[(3.14)]{BalabanRG1}, built from the localised functions of \cite[p.~273]{BalabanRG1} with the variables $s$ of \cite[p.~8]{BalabanRG2};
\item in \cite[(3.33)--(3.36), p.~277]{BalabanRG1} and \cite[(3.53), p.~280]{BalabanRG1}, at the background field and current of $\exp iw$, $w = \tau B + B'$ with $\tau \in [0, 1]$ and $B'$ small, where $B$ is built from the localised function of \cite[(1.8), p.~3]{BalabanRG2}; \cite[p.~276]{BalabanCRE} applies the same formulae to the boundary terms;
\item in the large-field steps, in the two localisations of \cite[(1.59)--(1.60), p.~372]{BalabanLFR2} and \cite[pp.~373--379]{BalabanLFR2}, at the configurations of the passes and at the points of the Taylor remainder of \cite[(3.34)]{BalabanRG1}, and in the terms $V(Y)$ and $V(X^\sim)$ of \cite[(1.70)--(1.72), pp.~377--379]{BalabanLFR2}, at the configurations \cite[(1.61)--(1.65), pp.~373--375]{BalabanLFR2}, which contain the gauge transformations ``given by the explicit formula (106) [12]'' \cite[p.~373]{BalabanLFR2}, that is \cite[(106), p.~33]{BalabanCMP98}.
\end{enumerate}
At real $(\mathbf U, \mathbf J)$ and real values of the auxiliary variables each of these arguments is a real point of the domain of the earlier term. Indeed, the localised operators are sums over walks whose terms are multiplied by the real numbers $s(\Delta)$ \cite[p.~3]{BalabanRG2}, so that the localised functions are real at real arguments, and so are those of \cite[(1.54), (1.57), (1.58), pp.~370--371]{BalabanLFR2}, built in the same way from the operators of \cite{BalabanPBG,BalabanVAR}. The gauge transformations of \cite[(3.26)--(3.27), p.~275]{BalabanRG1} are ``explicitly given analytic functions'' of $\mathbf U$ \cite[p.~276]{BalabanRG1}, $G$-valued at $G$-valued $\mathbf U$. In \cite[(106), p.~33]{BalabanCMP98} the transformation is a product of conjugations $R(X)Y = XYX^{-1}$ \cite[(56), p.~27]{BalabanCMP98} by contour products and of inverses of averages \cite[(105)]{BalabanCMP98}, which are exponentials of real averages of logarithms, so that it is $G$-valued at $G$-valued configurations. Every term read at these arguments is continuous in $h$ there, by (a) of $S_{j-1}$ for a term of an earlier step and by Step~3 for a term produced earlier in the step $j$.

\emph{Derivatives at real points.} Let $F$ satisfy (a) and let $p$ be a real point of $D_F$ at $h^0$. There is a complex polydisc $P \ni p$ contained in $D_F$ for all $h$ near $h^0$, because $D_F$ does not move with $h$ except for the bound of the boundary terms, which is not attained at $p$. On $P$ the term is bounded uniformly for $h$ near $h^0$ by the bound of its class, \cite[(1.18), p.~263]{BalabanRG1}, \cite[(2.31), p.~260]{BalabanCRE}, \cite[(2.42), p.~261]{BalabanCRE}, \cite[(1.70), p.~192]{BalabanLFR1}, or \cite[(1.11), p.~358, (1.47)--(1.48), p.~368]{BalabanLFR2}, with constants which do not depend on the history (Paper, Proposition~\ref{P-prop:construction-admissible}). By Montel's theorem every sequence $h_n \to h^0$ has a subsequence along which $F_{h_n}$ converges locally uniformly on $P$; every limit equals $F_{h^0}$ at the real points of $P$, which form a nonempty open subset of the real space, and therefore on $P$, by the identity theorem. So $F_h \to F_{h^0}$ locally uniformly on $P$, with all derivatives in the configuration, and finitely many such polydiscs cover a compact set of real points.

\emph{The auxiliary variables.} The localisations of \cite[(1.9)--(1.10), p.~4]{BalabanRG2} and \cite[(3.15), p.~273]{BalabanRG1} differentiate the integrand $\varphi_h(t_\square, s)$ at $t_\square = 0$ and $s \in [0, 1]^{Y_0}$, representing the derivatives by Cauchy integrals over circles \cite[p.~7]{BalabanRG2}. On the open polydisc $D$ of these circles, in which $\abs{s(\Delta)} < e^{\kappa_1}$, the field $\mathbf A$ stays strictly inside the class \cite[(3.14)]{BalabanRG1}, or \cite[(3.31), p.~276]{BalabanRG1}, so that $\varphi_h$ is analytic on $D$ and bounded there by \cite[(3.17), p.~273, (3.54), p.~280]{BalabanRG1}, uniformly in $h$. In case (ii), $\varphi_h$ is the remainder of \cite[(3.34)]{BalabanRG1}, an integral over $\tau \in [0, 1]$ of fifth derivatives, or a Taylor term of that formula; the function $w \mapsto \mathbf E^{(j)}(X, U_j(\square_0, \exp iw))$ is analytic and bounded on a connected open set of the $w$ above \cite[(3.53)--(3.54)]{BalabanRG1}, and it converges at the real points of that set by the preceding paragraph, so that by Vitali's theorem it converges uniformly with all derivatives on $\{ \tau B + B' \, | \, \tau \in [0, 1],\ \abs{B'} \leq \frac12\alpha_3 \}$. The transformations of \cite[Sect.~4]{BalabanRG1} act only on the Taylor terms at $B = 0$ (``such expressions only appear in the sum (4.6)'', \cite[p.~284]{BalabanRG1}), and by \cite[(4.2)--(4.3), p.~281]{BalabanRG1} and \cite[(4.9)--(4.15), pp.~283--284]{BalabanRG1} these are sums of derivatives of order at most four at $w = 0$ contracted with fields and cutoffs which do not depend on $h$. So in cases (i) and (ii) $\varphi_h$ converges at the real points of $D$, and by Vitali's theorem uniformly on the compact subsets of $D$, with all derivatives; the set $\{ t_\square = 0 \} \times [0, 1]^{Y_0}$ is compact in $D$ because $e^{\kappa_1} > 1$, and the sums over $Y_0$ converge uniformly in $h$ with \cite[(1.24)--(1.29), pp.~7--8]{BalabanRG2} as majorants. In case (iii) the expansions ``(6.9), (6.10) [I]'' which \cite[pp.~372--374]{BalabanLFR2} uses are \cite[(1.9)--(1.10), p.~4]{BalabanRG2}, obtained by the fundamental theorem of calculus in each $s(\Delta)$, so that a term is the finite alternating sum $\sum_{Z \subset \sigma}(-1)^{\abs{\sigma \setminus Z}}$ of the expression at $s = \mathbf 1_Z$, and no complex $s$ is used. At $s = \mathbf 1_Z$ the localised functions contain no coupling: the configurations of \cite[(1.50), p.~370]{BalabanLFR2} at $B = 0$ and of \cite[(1.53), (1.56)]{BalabanLFR2} are background fields of $G$-valued fields, the localised functions come from \cite[Sects.~D, E]{BalabanVAR} with the operators of \cite{BalabanPBG}, and the function of \cite[(1.51)]{BalabanLFR2} is read at the unscaled field $g_kB$. In the second localisation the expression at $s = \mathbf 1_Z$ is a term read at a real point of its domain, which converges by (a); in the first it is, as in \cite[(3.34)]{BalabanRG1}, a sum of Taylor terms of order at most four at a real configuration and an integral over $\tau$ of fifth derivatives at the real configurations between two members, contracted with real fields, and these converge by the preceding paragraph. The sums and the resummations converge uniformly in $h$, with the bounds of \cite[(1.24)--(1.29), pp.~7--8]{BalabanRG2}, those named on \cite[p.~372]{BalabanLFR2}, and the factors \cite[(1.66), p.~376]{BalabanLFR2} as majorants.

\emph{Step 3: the terms of the step at real points.} Each term of the step $j$ is continuous in $h$ at every real point:
\begin{itemize}
\item \emph{Fluctuation integrals.} The integrands are the densities of Step~1 at real unscaled fields, restricted by fixed characteristic functions to fields which lie in the real part of the domains: the thresholds $\varepsilon_j$ and $\vartheta_j$ are of order $(\log\hat g_j^{-2})^{p_0-q_i}$ relative to $\alpha_{i,j}$, and $q_0 = q_1 > p_0 + 3r$ (Paper, Appendix~\ref{P-app:constants}, the stage $0$). The densities are continuous in $h$, by the continuity of $\beta_i$, $i \leq j$, in the coefficient function \eqref{P-eq:history-coupling}, by Step~2, and by $S_{j-1}$, and they are bounded uniformly for $h$ near $h^0$; dominated convergence gives the continuity of each integral.
\item \emph{Small-field terms.} The term $\mathbf E^{(j+1)}(X)$ is the series \cite[(2.13), p.~14]{BalabanRG2} in the activities \cite[(2.14), p.~15]{BalabanRG2}, which are integrals over $s(\Delta), t(Y) \in [0, 1]$ of Cauchy integrals of Gaussian integrals in which ``the operators in it are not symmetric, and the second measure is complex'' \cite[p.~15]{BalabanRG2}. In the unscaled fields, at a real point and at fixed values on the circles, $h$ enters the integrand only through $g_j$ in the Gaussian densities $\exp(-Q/(2g_j^2))$, with a form $Q$ which does not depend on $h$, and through the potentials, which converge by Step~2 and $S_{j-1}$. \cite[(2.15)--(2.25), pp.~15--17]{BalabanRG2} bound the modulus of the integrand by a constant times a positive Gaussian density in which $g_j^{-2}$ multiplies a nonnegative form, and for $\frac12g^0_j \leq g_j \leq 2g^0_j$ this is at most $4^n$ times its value at $2g^0_j$, $n$ the number of the variables. Dominated convergence gives the continuity of the activities, and the series converges uniformly in $h$ by \cite[(2.38), p.~20]{BalabanRG2} and ``the standard way'' of \cite[p.~20]{BalabanRG2}.
\item \emph{The $\mathbf R$-operation.} The quotients of \cite[(0.3), (0.5), p.~176]{BalabanLFR1}, with the densities and normalisations of \cite[(1.100), p.~201]{BalabanLFR1}, have positive denominators \cite[p.~176]{BalabanLFR1}, which are continuous in $h$ by the first item; for the soft cutoff of Definition~\ref{def:soft-functions} the denominators are positive by Proposition~\ref{prop:soft-tests}(d). The new terms $\mathbf R'^{(k)}(X)$ of \cite[(1.98)--(1.100), p.~390]{BalabanLFR2} are ``given by the convergent series (7.13) [I]'' \cite[p.~390]{BalabanLFR2}, which is \cite[(2.13)]{BalabanRG2}, in the activities \cite[(1.91), p.~388]{BalabanLFR2}, which are built from the operation $T'_k$ of \cite[(1.71), p.~379]{BalabanLFR2}. In $T'_k$ ``only the quadratic form in $B$ and the term $V(X^\sim)$ depend on $(\mathbf U, \mathbf J)$'' \cite[p.~379]{BalabanLFR2}, the integration variables are real at a real point, the Haar density of the bond, read at the unscaled field, is bounded on the small fields which the characteristic function admits, and the real part of the exponent is bounded above uniformly for $h$ near $h^0$: the quadratic form and the Wilson action are nonnegative, $V(X^\sim)$ is bounded by \cite[(1.69), p.~377]{BalabanLFR2}, and the counterterms and normalisations by $O(\log g_j^{-2})$ per site \cite[p.~380]{BalabanLFR2}. With the Gaussian integrals dominated as in the previous item and the gauge fields in compact groups, dominated convergence gives the continuity of $T'_k$, and the series converge uniformly with \cite[(1.92)--(1.97), (1.99)--(1.100), pp.~388--390]{BalabanLFR2} as majorants.
\item \emph{The terms $\mathbf C_k$.} They are ``the sum of all these terms, including the term $I(U_0)$ from the denominator'' \cite[p.~369]{BalabanLFR2}, obtained ``by expanding some Wilson actions'' \cite[p.~369]{BalabanLFR2}. A Wilson action is linear in its coefficient function, and the expanded configurations contain no coupling (Step~1), so that at a real point these terms are finite sums of plaquette functions at $G$-valued configurations with coefficients which are continuous in $h$. The term $I(U_0)$ of \cite[(1.10), p.~358]{BalabanLFR2} is continuous by dominated convergence, since its integral $I_\Lambda$ is over a product of copies of $G$ of a nonnegative integrand bounded by one, and is positive by the lower bound of \cite[p.~358]{BalabanLFR2}.
\end{itemize}
The step produces its terms in the order: the $N_j$ preliminary integrations of \cite[pp.~188--189]{BalabanLFR1}, each a fluctuation integral ``of the same type as the integral studied in Sect.~3 [III]'' \cite[p.~189]{BalabanLFR1} read as in cases (i) and (ii); the terms $\mathbf C_j$; the first localisation and its repetitions, and the second localisation; the $\mathbf R$-operation; and the renormalisation transformation of \cite[Sect.~3]{BalabanCRE}. Each operation reads only terms of the earlier steps and terms produced before it, so that (a) holds for every term produced by the step $j$.

\emph{Step 4: polydiscs.} Let $F$ be a regular term or an $\mathbf R$-term of the step $j$, $B^0$ admissible, and $f_h(B) \coloneqq F(X, u_j(B); h)$ on $P_j(B^0)$. At real $B$ the point $u_j(B)$ is a real point of $D_F$, so that $f_h(B) \to f_{h^0}(B)$ by Step~3. On $P_j(B^0)$ the functions $f_h$ are analytic and bounded uniformly for $h$ near $h^0$, by \cite[(1.18), p.~263]{BalabanRG1} or \cite[(2.31), p.~260]{BalabanCRE}. As in Step~2, Montel's theorem and the identity theorem on the connected set $P_j(B^0)$ give $f_h \to f_{h^0}$ locally uniformly, which is (b).

\emph{Step 5: the coefficient.} The coefficient is $\Lambda$ of the regular part on $\mathbb Z^4$, obtained from the tori $T^{(j+1)} \nearrow \mathbb Z^4$ (``This limit exists by the localized representation (1.7)'', \cite[p.~264, after (1.21)]{BalabanRG1}), and the finite-volume values are continuous in $h$ by Steps~1--4 and (b) at $B^0 = 0$. It remains to show that the limit is uniform for $h$ near $h^0$. Within one step the localised terms are built from the random walk expansions of the propagators \cite[(1.6)--(1.7), p.~3]{BalabanRG2}, from \cite[(3.105)--(3.108), pp.~414--416]{BalabanPBG}, whose walks $\omega$ are products of local operators, each depending on the configuration on its own domain, with the bound \cite[(1.7), p.~3]{BalabanRG2}, which decays exponentially in the length $\abs{\omega}$ of the walk. For a fixed $X$ the torus of side $\lambda_T$ enters only through the walks of length at least a fixed multiple of $\lambda_T$, so that the term on the torus differs from the term on $\mathbb Z^4$, to which \cite[p.~281]{BalabanCRE} extends the sums, by at most $C_je^{-c_{\mathrm{tor}}\lambda_T}$ times its bound. The inputs of the step come from the torus as well, and an induction over the steps $0, \dots, j$, with Cauchy's estimates on the domains of \cite[p.~259]{BalabanCRE}, carries this bound to every term with $\diam X \leq \lambda_T/3$; the constant $C_j$ may grow with $j$, but only $k + 1$ steps enter. On the other hand the proof of Lemma~\ref{lem:polarisation-bounded} gives the majorant
\begin{equation}\label{eq:polarisation-majorant}
    \abs{\Pi_X(x; h)}\abs{x}^2 \leq m(X, x) \coloneqq CE_0r^{-2}e^{-\kappa d_{j+1}(X)}e^{-\delta_0(\dist(0, X) + \dist(x, X))/2}\abs{x}^2\;,
\end{equation}
uniform in $h$ and in the torus, with $\sum_{X,x}m(X, x) < \infty$ on $\mathbb Z^4$. Given $\varepsilon > 0$, a finite set of pairs $(X, x)$ near $0$ carries all but $\varepsilon$ of this sum; for $\lambda_T$ large these pairs are pairs of the torus with the same distances, and on them the torus terms converge uniformly in $h$. A pair of the torus which is not of this kind has $\diam X$, $\dist(0, X)$, or $\abs{x}$ at least $\lambda_T/8$, or $d_{j+1}(X) \geq \lambda_T/(4M)$, and then one of the factors of $m(X, x)$ is at most $e^{-c\lambda_T}$, by $\diam X \leq M(2d_{j+1}(X) + 4)$ (Lemma~\ref{lem:cube-count}); the sum of the majorants of these pairs is at most a polynomial in $\lambda_T$ times $e^{-c\lambda_T}$. Hence $\sum_x\Pi_{j+1}(x; h)x_\mu x_\nu$ converges uniformly for $h$ near $h^0$, and $\beta_{j+1}$, a uniform limit of continuous functions, is continuous at $h^0$. Since $E^{(j+1)}$ contains $-\beta_{j+1}A(\phi_{j+1}, \cdot)$ \cite[(2.25), p.~259]{BalabanCRE}, it satisfies (a) and (b) as well, which completes $S_j$. The claim $S_k$ contains the continuity of $\beta_{k+1}$ on $H_k$.
\end{proof}

\section{The Coefficients at Zero Coupling}\label{sec:zero-coupling}

Paper, Theorem~\ref{P-thm:one-loop-window}, states that the coefficients $b_j = \Lambda(\log Z^{(j-1)})$ of \eqref{P-eq:beta-split} add up, over any number of consecutive steps, to $2c\ln L$ per step up to a bounded error. Ba{\l}aban states the corresponding result for the running coupling as \cite[Theorem~2, p.~259]{BalabanRG1}, for $G = SU(2)$, and writes: ``A proof of this theorem, based on perturbative calculations, will be given in a separate paper''. This section gives the proof, for a simple group. A single factor $Z^{(j)}$ is the Gaussian normalisation of one step, and its polarisation has no simple form, because the quadratic form of the step contains the averaging of all the earlier steps. The difficulty disappears in the sum: we show that $\sum_{j<k}\log Z^{(j)}$ is, up to terms which do not contribute to the coefficient, the difference $F_k - F_0$ of two functionals, each of which is the logarithm of the determinant of a gauge-invariant form, of the action seen at the level $k$ and at the level $0$. The polarisation of $F_k$, measured on the lattice of the level $k$, is that of a theory with a cutoff at the unit scale on both sides and is bounded uniformly in $k$, while that of $F_0$ is the one-loop polarisation of the bare lattice theory seen from the lattice of the level $k$, and grows like $2ck\ln L$ times the action of the background. The terms left over come from the current $J$ of \cite[(1.8), p.~261]{BalabanRG1}, and gauge covariance makes their polarisation vanish to the order which matters.

\subsection{An Exact Decomposition}\label{sec:exact-decomposition}

We first prove the identity, in the setting of finite-dimensional linear algebra in which it is exact.

\emph{The setting.} Let $X$, $Y$, $Z$, and $\bar Y$ be finite-dimensional real inner-product spaces, and let
\begin{equation}\label{eq:linear-setting}
    D \colon Y \to X\;,
    \qquad
    h \colon X \to X\;,
    \qquad
    Q \colon X \to Z\;,
    \qquad
    P \colon Y \to \bar Y\;,
    \qquad
    \bar D \colon \bar Y \to Z
\end{equation}
be linear maps such that $D$ and $\bar D$ are injective, $h$ is symmetric and nonnegative with $hD = 0$ and positive definite on $(\operatorname{ran}D)^\perp$, $Q$ is onto, $PP^{\mathrm T} = 1$, and
\begin{equation}\label{eq:covariance-of-average}
    QD = \bar DP\;.
\end{equation}
In the application $X$ is a space of fine Lie-algebra-valued $1$-forms, $Y$ of fine $0$-forms orthogonal to $\ker D$, $D = D_U$ the linearised action of the gauge group at a background $U$, $Q$ the linearised $k$-fold average with values in the coarse $1$-forms $Z$ of the level $k$, $P$ the restriction to the sites of the level $k$, which are fine sites because $L$ is odd \cite[p.~251]{BalabanRG1}, and $\bar D$ the coarse gauge action. The relation \eqref{eq:covariance-of-average} is the linearisation of the gauge covariance of the average: the averages of \cite[(0.4), (0.12), pp.~253--254]{BalabanRG1} are built from contours based at the sites of the coarse lattice, which transform as $U(y, x) \mapsto u(y)U(y, x)u(x)^{-1}$ \cite[(0.6)]{BalabanRG1}, so that the average of $U^u$ is $u(c_-)\bar U(c)u(c_+)^{-1}$, and the linearisation of this, iterated over the $k$ steps, is \eqref{eq:covariance-of-average}. Put $Y_0 \coloneqq \ker P$, the gauge transformations which are trivial at the coarse sites, $D_0 \coloneqq D|_{Y_0}$, and $W \coloneqq \ker Q \ominus D(Y_0)$, where $D(Y_0) \subset \ker Q$ by \eqref{eq:covariance-of-average}.

\begin{definition}[Forms and Functionals of the Levels]\label{def:level-functionals}
In the setting \eqref{eq:linear-setting}, the \emph{one-loop factor} of the $k$-fold average is
\begin{equation}\label{eq:one-loop-factor}
    Z_k[h] \coloneqq \det(QQ^{\mathrm T})^{-1/2}\det(h|_W)^{-1/2}\det(D_0^{\mathrm T}D_0)^{1/2}\;.
\end{equation}
For a level $j$, with the average $Q_j$, the restriction $P_j$, and the coarse gauge action $\bar D_j$, the form $\Delta_j$ on $Z_j$ is
\begin{equation}\label{eq:level-form}
    \langle z, \Delta_jz \rangle \coloneqq \min\bigl\{ \langle A, hA \rangle \, \big| \, Q_jA \in z + \operatorname{ran}\bar D_j \bigr\}\;,
\end{equation}
which depends on $z$ modulo $\operatorname{ran}\bar D_j$, and
\begin{equation}\label{eq:level-functional}
    F_j \coloneqq \tfrac12\log\det{}_\perp(\Delta_j) - \tfrac12\log\det(\bar D_j^{\mathrm T}\bar D_j)\;,
\end{equation}
with $\det_\perp$ the determinant on $Z_j \ominus \operatorname{ran}\bar D_j$. At the level $0$, where $Q_0 = 1$ and $\bar D_0 = D$, $\Delta_0 = h$ because $hD = 0$, and $F_0 = \frac12\log\det_\perp(h) - \frac12\log\det(D^{\mathrm T}D)$.
\end{definition}

The minimum in \eqref{eq:level-form} is attained, because a convex quadratic function which is bounded below on a finite-dimensional affine space attains its infimum. The factor \eqref{eq:one-loop-factor} is the constrained Gaussian integral of $e^{-\frac12\langle A, hA \rangle}\delta(QA)$ over a slice of the residual gauge orbits, in the axial gauge, in which the Faddeev--Popov determinant is one: the map $(A_s, \omega) \mapsto A_s + D\omega$ from the slice times $Y_0$ has the block-triangular Jacobian $\det D_{\mathrm{tree}}$, where $D_{\mathrm{tree}}$ is $\omega \mapsto D\omega$ on the bonds of the trees of the axial gauge, and ordering these bonds from the coarse sites outwards makes $D_{\mathrm{tree}}$ triangular with orthogonal diagonal blocks, so that $\abs{\det D_{\mathrm{tree}}} = 1$. The integral of the $\delta$-function is $\det(QQ^{\mathrm T})^{-1/2}$ times the integral over $\ker Q$, and the splitting $\ker Q = W \oplus D(Y_0)$, with the orbit measure $\det(D_0^{\mathrm T}D_0)^{1/2}\,\dd\omega$, gives \eqref{eq:one-loop-factor} up to a constant which does not depend on the background.

\begin{proposition}[The Exact Decomposition]\label{prop:exact-decomposition}
In the setting \eqref{eq:linear-setting}, $\log Z_k[h] = F_k - F_0$.
\end{proposition}

\begin{proof}
Let $K \coloneqq \ker Q$ and $\Delta_F \coloneqq h + DD^{\mathrm T}$. With respect to $X = \operatorname{ran}D \oplus (\operatorname{ran}D)^\perp$, which $h$ preserves because it is symmetric with $hD = 0$, $\Delta_F$ is $DD^{\mathrm T}$ on $\operatorname{ran}D$, positive there because $D$ is injective, and $h$ on the complement, positive there by hypothesis; so $\Delta_F$ is invertible, with
\begin{equation}\label{eq:feynman-determinant}
    \det\Delta_F = \det(D^{\mathrm T}D)\det{}_\perp(h)\;,
\end{equation}
the nonzero eigenvalues of $DD^{\mathrm T}$ and $D^{\mathrm T}D$ being the same. Moreover $\Delta_FD = D(D^{\mathrm T}D)$, hence $\Delta_F^{-1}D = D(D^{\mathrm T}D)^{-1}$. We write $S \coloneqq Q\Delta_F^{-1}Q^{\mathrm T}$ and $R \coloneqq P(D^{\mathrm T}D)^{-1}P^{\mathrm T}$.

\emph{The form with the residual gauge directions.} Since $PP^{\mathrm T} = 1$, $P^{\mathrm T}P$ is the orthogonal projection onto $(\ker P)^\perp$, so that $D_0D_0^{\mathrm T} = D(1 - P^{\mathrm T}P)D^{\mathrm T}$. Let $h_F$ be the compression of $h + D_0D_0^{\mathrm T}$ to $K$. For $w \in W$ and $\omega \in Y_0$ we have $hD\omega = 0$, $\langle D\omega, hw \rangle = 0$, and $D_0^{\mathrm T}w = 0$, the last because $w \perp D(Y_0)$. So $h_F$ is the direct sum of $h|_W$ and of $D_0D_0^{\mathrm T}$ on $D(Y_0)$, whose determinant is $\det(D_0^{\mathrm T}D_0)$, and $\det h_F = \det(h|_W)\det(D_0^{\mathrm T}D_0)$. With \eqref{eq:one-loop-factor},
\begin{equation}\label{eq:one-loop-feynman}
    \log Z_k[h] = -\tfrac12\log\det(QQ^{\mathrm T}) - \tfrac12\log\det h_F + \log\det(D_0^{\mathrm T}D_0)\;.
\end{equation}

\emph{The ghost determinant.} With respect to $Y = Y_0 \oplus \operatorname{ran}P^{\mathrm T}$, the Schur complement of the block $D_0^{\mathrm T}D_0$ of $D^{\mathrm T}D$ is the inverse of the block of $(D^{\mathrm T}D)^{-1}$ on $\operatorname{ran}P^{\mathrm T}$, which $P^{\mathrm T}$ identifies isometrically with $R$. Hence
\begin{equation}\label{eq:ghost-schur}
    \det(D_0^{\mathrm T}D_0) = \det(D^{\mathrm T}D)\det R\;.
\end{equation}

\emph{The constraint.} We have $h + D_0D_0^{\mathrm T} = \Delta_F - DP^{\mathrm T}PD^{\mathrm T}$. The compression of $\Delta_F$ to $K$ has the determinant $\det\Delta_F\det S/\det(QQ^{\mathrm T})$, and its inverse, extended by zero, is $C \coloneqq \Delta_F^{-1} - \Delta_F^{-1}Q^{\mathrm T}S^{-1}Q\Delta_F^{-1}$; these are the two formulae of the constrained Gaussian integral. The determinant of a rank perturbation, $\det(1 - UV) = \det(1 - VU)$, gives
\begin{equation}\label{eq:determinant-lemma}
    \det h_F = \frac{\det\Delta_F\det S}{\det(QQ^{\mathrm T})}\det{}_{\bar Y}(1 - PD^{\mathrm T}CDP^{\mathrm T})\;.
\end{equation}
By $\Delta_F^{-1}D = D(D^{\mathrm T}D)^{-1}$ and \eqref{eq:covariance-of-average}, $D^{\mathrm T}\Delta_F^{-1}D = 1$ and $Q\Delta_F^{-1}D = \bar DP(D^{\mathrm T}D)^{-1}$, so that
\begin{equation}\label{eq:ghost-constraint}
    1 - PD^{\mathrm T}CDP^{\mathrm T} = R\bar D^{\mathrm T}S^{-1}\bar DR\;,
    \qquad
    \det{}_{\bar Y}(1 - PD^{\mathrm T}CDP^{\mathrm T}) = (\det R)^2\det(\bar D^{\mathrm T}S^{-1}\bar D)\;.
\end{equation}

\emph{Collecting.} Inserting \eqref{eq:ghost-schur}, \eqref{eq:determinant-lemma}, and \eqref{eq:ghost-constraint} into \eqref{eq:one-loop-feynman}, the factors $\det(QQ^{\mathrm T})$ cancel, and so do the factors $\det R$, which enter with $-\frac12 \cdot 2 + 1 = 0$:
\begin{equation}\label{eq:two-brackets}
    \log Z_k[h] = \Bigl[-\tfrac12\log\det\Delta_F + \log\det(D^{\mathrm T}D)\Bigr] + \Bigl[-\tfrac12\log\det S - \tfrac12\log\det(\bar D^{\mathrm T}S^{-1}\bar D)\Bigr]\;.
\end{equation}
By \eqref{eq:feynman-determinant} the first bracket is $-\frac12\log\det_\perp(h) + \frac12\log\det(D^{\mathrm T}D) = -F_0$. For the second, let $M \coloneqq S^{-1}$, so that $\langle z, Mz \rangle = \min\{ \langle A, \Delta_FA \rangle \, | \, QA = z \}$. With respect to $Z = \operatorname{ran}\bar D \oplus (\operatorname{ran}\bar D)^\perp$, $\det M$ is the product of the determinant of the block $M_\parallel$ on $\operatorname{ran}\bar D$, which is $\det(\bar D^{\mathrm T}M\bar D)/\det(\bar D^{\mathrm T}\bar D)$, and of the determinant of the Schur complement $M/M_\parallel$. So the second bracket is $\frac12\log\det(M/M_\parallel) - \frac12\log\det(\bar D^{\mathrm T}\bar D)$. The Schur complement is the form
\begin{equation}\label{eq:schur-level}
    \langle z, (M/M_\parallel)z \rangle = \min\bigl\{ \langle A, (h + DD^{\mathrm T})A \rangle \, \big| \, QA \in z + \operatorname{ran}\bar D \bigr\}\;,
\end{equation}
and the set of the $A$ is invariant under $A \mapsto A + D\omega$, because $QD\omega = \bar DP\omega$. As $h$ does not change under this map, and $\omega = -(D^{\mathrm T}D)^{-1}D^{\mathrm T}A$ makes $D^{\mathrm T}(A + D\omega) = 0$, the minimum \eqref{eq:schur-level} is $\langle z, \Delta_kz \rangle$, and the second bracket is $F_k$.
\end{proof}

At $k = 0$ the constraint fixes everything and both sides vanish. The steps of the construction average one level at a time, and the proposition applies to each of them.

\begin{corollary}[One Step]\label{cor:one-step}
Let $\tilde Q$ be the linearised average from the level $k$ to the level $k + 1$ and $\tilde P$ the restriction from the sites of the level $k$ to those of the level $k + 1$, so that $Q_{k+1} = \tilde QQ_k$ and $P_{k+1} = \tilde PP_k$. Let $Z^{(k)}[h]$ be the factor \eqref{eq:one-loop-factor} of the setting in which $X$, $h$, $Y$, $D$, $Q$, $P$, $\bar D$ are replaced by $Z_k$, $\Delta_k$, $\bar Y_k$, $\bar D_k$, $\tilde Q$, $\tilde P$, $\bar D_{k+1}$. Then $\log Z^{(k)}[h] = F_{k+1} - F_k$, and $\sum_{j<k}\log Z^{(j)}[h] = \log Z_k[h]$.
\end{corollary}

\begin{proof}
The replacements satisfy the hypotheses of the setting. Indeed, $\Delta_k$ vanishes on $\operatorname{ran}\bar D_k$ and is positive on its orthogonal complement, by Lemma~\ref{lem:full-space-form} below, whose proof does not use the corollary, so that $\Delta_k\bar D_k = 0$; $\tilde Q$ is onto because $Q_{k+1} = \tilde QQ_k$ is; $\tilde P\tilde P^{\mathrm T} = 1$; $\bar D_k$ and $\bar D_{k+1}$ are injective; and $\tilde Q\bar D_kP_k = \tilde QQ_kD = Q_{k+1}D = \bar D_{k+1}\tilde PP_k$, with $P_k$ onto, gives $\tilde Q\bar D_k = \bar D_{k+1}\tilde P$. The functional of the level $0$ of this setting is $F_k$. Its form of the level $1$ is $\min\{ \langle z', \Delta_kz' \rangle \, | \, \tilde Qz' \in z + \operatorname{ran}\bar D_{k+1} \} = \min\{ \langle A, hA \rangle \, | \, Q_{k+1}A \in z + \operatorname{ran}\bar D_{k+1} \}$, since $\tilde Q(\operatorname{ran}\bar D_k) = \bar D_{k+1}\tilde P(\bar Y_k) \subset \operatorname{ran}\bar D_{k+1}$, and $z' = Q_kA$ gives the converse inequality; this is $\langle z, \Delta_{k+1}z \rangle$, and the functional of the level $1$ is $F_{k+1}$. Proposition~\ref{prop:exact-decomposition} with one step gives the first identity, and the sum over $j < k$, compared with the proposition, gives the second.
\end{proof}

\begin{lemma}[The Form on the Whole Space]\label{lem:full-space-form}
The form $\Delta_k$ vanishes on $\operatorname{ran}\bar D_k$ and is positive definite on its orthogonal complement. Consequently, with $\mathcal D_k \coloneqq \Delta_k + \bar D_k\bar D_k^{\mathrm T}$,
\begin{equation}\label{eq:level-functional-full}
    F_k = \tfrac12\log\det\mathcal D_k - \log\det(\bar D_k^{\mathrm T}\bar D_k)\;.
\end{equation}
\end{lemma}

\begin{proof}
The form $\Delta_k \geq 0$ depends on $z$ only modulo $\operatorname{ran}\bar D_k$, so by polarisation it vanishes on $\operatorname{ran}\bar D_k$ and maps the orthogonal complement into itself. If $\langle z, \Delta_kz \rangle = 0$, a minimiser $A$ has $\langle A, hA \rangle = 0$, hence $A = D\omega$, because $h$ is positive definite on $(\operatorname{ran}D)^\perp$ and $hD = 0$, and then $z \in Q_kD\omega + \operatorname{ran}\bar D_k = \bar D_kP_k\omega + \operatorname{ran}\bar D_k = \operatorname{ran}\bar D_k$. Since $\bar D_k\bar D_k^{\mathrm T}$ vanishes on $(\operatorname{ran}\bar D_k)^\perp$ and has the determinant $\det(\bar D_k^{\mathrm T}\bar D_k)$ on $\operatorname{ran}\bar D_k$, $\log\det_\perp\Delta_k = \log\det\mathcal D_k - \log\det(\bar D_k^{\mathrm T}\bar D_k)$, which with \eqref{eq:level-functional} gives \eqref{eq:level-functional-full}.
\end{proof}

\bibliographystyle{amsplain}
\bibliography{refs}
\end{document}
